\documentclass[11pt,a4paper]{report}

\usepackage[utf8]{inputenc}
\usepackage[T1]{fontenc}
\usepackage[english]{babel}
\usepackage{lmodern}
\usepackage[margin=2.5cm]{geometry}
\usepackage{graphicx}
\usepackage{hyperref}
\usepackage{xcolor}
\usepackage{tikz}
\usetikzlibrary{trees,positioning}
\usepackage{enumitem}
\usepackage{url}
\usepackage{booktabs}
\usepackage{amsmath}

\hypersetup{
  colorlinks=true,
  linkcolor=black,
  urlcolor=blue,
  citecolor=black
}

\title{Towards an Automated Source Code Remediation Approach \\
       for the Security and Privacy of Third-Party Mobile Applications}
\author{Master's Thesis \\
        Master in Sciences and Technologies \\
        Major: Computer Science, Track: Software Engineering}
\date{}

\begin{document}

% =====================================================================
% TITLE PAGE
% =====================================================================
\begin{titlepage}
\begin{center}
    % Logo/En-tête (image complète)
    \includegraphics[width=\textwidth]{logo.png}\\[2cm]
    
    % Titre du document
    {\LARGE\bfseries Graduation thesis}\\[0.8cm]
    {\large For the obtention of the State Computer Science Engineer Diploma}\\[0.5cm]
    {\large\bfseries Option: Computer Systems (SQ)}\\[1cm]
    
    \rule{\linewidth}{0.5pt}\\[0.5cm]
    % Titre du projet
    {\large\bfseries Topic:}
    \vspace{0.3cm}
    
    {\LARGE Towards an Approach for Software Security \\[0.3cm]Vulnerability Identification and Remediation}\\[0.5cm]
    \rule{\linewidth}{0.5pt}\\[1.5cm]
    
    % Auteurs et encadrants
    \noindent
    \begin{minipage}{0.45\textwidth}
        \begin{flushleft}\large
            \textit{Made by :}\\
            - \textsc{Nacef} Salah Eddine\\
            - \textsc{Nouali} Asmaa
        \end{flushleft}
    \end{minipage}%
    \hfill
    \begin{minipage}{0.5\textwidth}
        \begin{flushright}\large
            \textit{Supervised by :}\\
            - Mr. \textsc{Seriai} Abdelhak-Djamel (LIRMM)\\
            - Mrs. \textsc{Zellagui} Soumia (ESI)
        \end{flushright}
    \end{minipage}\\[1cm]
    
    \vfill
    {\large Academic Year: 2025/2026}
\end{center}
\end{titlepage}

%=====================================================================
% DEDICATION
% =====================================================================
\chapter*{Dedication}
\addcontentsline{toc}{chapter}{Dedication}

%=====================================================================
% ACKNOWLEDGMENTS
% =====================================================================
\chapter*{Acknowledgments}
\addcontentsline{toc}{chapter}{Acknowledgments}

% =====================================================================
% ABSTRACT
% =====================================================================
\chapter*{Abstract}
\addcontentsline{toc}{chapter}{Abstract}

The widespread integration of software systems into daily activities has significantly increased the exposure of sensitive data to security risks. As applications become more complex and interconnected, they present a growing number of potential entry points for attackers. Ensuring the security of these systems is therefore a major concern, requiring effective mechanisms to identify and address vulnerabilities.

The process of securing software can be seen as a continuous confrontation between attackers, who aim to exploit weaknesses, and defenders, who seek to prevent or mitigate such exploits. In this context, vulnerability detection and remediation play a central role in reducing the attack surface and improving overall system resilience. While numerous techniques exist to uncover vulnerabilities, the challenge lies in efficiently correcting them without negatively impacting system functionality.

This work focuses on the analysis and improvement of approaches for detecting and remediating software vulnerabilities. We first identify a set of common vulnerabilities that may affect software systems, then propose corresponding remediation strategies designed to preserve application behavior. Finally, we evaluate our approach through its implementation within an automated framework, demonstrating its effectiveness in enhancing software security.

\medskip
\noindent\textbf{Keywords:} software security, vulnerability detection, automated remediation, automated program repair, static analysis, dynamic analysis, machine learning, large language models, systematic mapping study, secure software engineering.

%=====================================================================
% TABLE OF CONTENTS
% =====================================================================
\tableofcontents

% =====================================================================
% GENERAL INTRODUCTION
% =====================================================================
\section*{General Introduction}
\label{chap:introduction}

\addcontentsline{toc}{chapter}{General Introduction}


\subsection*{Context \& Problem Statement}
With the increasing dependence on software systems in almost every aspect of modern life, ensuring their security has become a critical concern. Applications today manage sensitive data, support essential services, and operate in highly interconnected environments. This growing complexity makes them particularly susceptible to a wide range of vulnerabilities that can be exploited by malicious actors.

Software vulnerabilities arise from design flaws, implementation errors, or misconfigurations within a system. These weaknesses can compromise confidentiality, integrity, and availability, leading to serious consequences such as data breaches, financial losses, or system failures. As software continues to evolve rapidly, the number of potential vulnerabilities increases, making their management a challenging task.

To address this issue, significant research efforts have focused on vulnerability detection. Various techniques have been developed, including static analysis, dynamic testing, and hybrid approaches, each offering different advantages in terms of accuracy, scalability, and coverage. More recently, intelligent methods leveraging machine learning have been introduced to enhance detection capabilities and reduce false positives.

However, identifying vulnerabilities is only part of the solution. Effective remediation is essential to ensure that detected issues are properly addressed. This process may involve applying patches, modifying source code, improving system configurations, or adopting secure development practices. The challenge lies in selecting appropriate remediation strategies that are both efficient and sustainable, without introducing new weaknesses.

\subsection*{Objectives}

This work pursues three main objectives. First, we aim to establish a comprehensive understanding of the existing techniques used to detect software vulnerabilities, covering static, dynamic, hybrid, and learning-based approaches, and to assess their respective strengths and limitations. Second, we investigate current automated remediation strategies, analyzing how Automated Program Repair techniques can be applied to address security flaws while preserving application behavior. Third, building on these foundations, we propose a graph-based framework that models privilege-related flows in web applications, combining graph neural network models with explainable reporting to identify access-control vulnerabilities, specifically missing authentication, missing authorization, and Insecure Direct Object Reference (IDOR). IDOR is a vulnerability that occurs when an application exposes internal object identifiers (such as database keys or file paths) directly to users without verifying that the requester is authorized to access the corresponding resource. The overall aim is to improve the effectiveness and integration of detection and remediation processes within a unified, automated pipeline.

\subsection*{Report Organization}
The remainder of this document is structured as follows:

\begin{itemize}[leftmargin=*]
\item \textbf{Part~I, State of the Art} comprises two chapters. Chapter~1 establishes the theoretical foundations by presenting a taxonomy of the research domain: the landscape of software security, the principal families of detection and remediation techniques, and the methodology of systematic mapping studies. Chapter~2 reviews the related work, existing secondary studies, and the key tools and benchmarks, and identifies the problems and research questions that motivate our contribution.

\item \textbf{Part~II, Contribution} comprises three chapters. Chapter~3 formalizes the problem of authorization vulnerability detection, positions our approach against existing paradigms, and presents the design of the Program Flow Graph (PFG) framework including the GNN architecture and structural validator. Chapter~4 details the implementation: framework-specific extractors, dataset construction from 75~real-world CVEs, training pipeline, and demonstration system. Chapter~5 evaluates the framework on the real-world benchmark, presents a failure-mode analysis, and discusses interpretability, scalability, and threats to validity.
\end{itemize}


% ============================================================
% PART II — CONTRIBUTION
% ============================================================

\part{Contribution}

\chapter{PFG: A Graph-Based Framework for Authorization Vulnerability Detection}
\label{chap:framework}

\section{Introduction}
\label{sec:framework-intro}

The state of the art examined in Part~I reveals two dominant paradigms for learning-based vulnerability detection in source code. The first family of approaches adopts Natural Language Processing techniques---most notably pre-trained code language models such as CodeBERT~\cite{feng2020codebert}, VulBERTa~\cite{hanif2022vulberta}, and LineVul~\cite{fu2022linevul}---that treat source code as token sequences and learn dense, 768-dimensional embeddings capturing generic code semantics. While these methods achieve competitive performance on benchmark datasets for injection and memory-safety vulnerabilities, they operate at the \emph{lexical} level: they detect surface patterns in code tokens rather than structural relationships between architectural components. For authorization vulnerabilities, this limitation is fundamental. Whether an endpoint is vulnerable depends not on how its code \emph{reads}, but on what guards---authentication middleware, role checks, ownership predicates---stand between it and the data access layer. A token-level model cannot reliably answer the question: \emph{does this endpoint have an authorization check before it queries the database?}

The second paradigm employs graph-based code representations, typically Code Property Graphs (CPGs) that unify the Abstract Syntax Tree (AST), Control Flow Graph (CFG), and Program Dependence Graph (PDG) into a single heterogeneous structure~\cite{yamaguchi2014modeling}. Graph Neural Networks (GNNs) are then trained over these representations, and post-hoc interpretability tools such as GNNExplainer~\cite{ying2019gnnexplainer} are applied to identify which subgraph was most influential for a given prediction. Although this approach captures structural code dependencies, it introduces considerable noise: a single function may yield thousands of AST and CFG nodes, the vast majority of which---assignments, literals, control-flow joins---carry no security-relevant information. GNNExplainer is necessary precisely \emph{because} the graph is opaque; it contains too many irrelevant nodes for direct human interpretation. Moreover, the explanation it produces (``these nodes contributed most to the prediction'') still requires additional interpretation to translate into actionable developer guidance.

\medskip

In this chapter, we present an approach that departs from both paradigms. We observe that authorization vulnerabilities are fundamentally \emph{architectural}: they arise from the absence of specific checks along the request-processing chain, not from syntactic patterns within individual functions. This observation motivates the design of the \textbf{Program Flow Graph~(PFG)}---a purpose-built, domain-specific graph representation that models exclusively the authorization-relevant structure of a web application. The PFG comprises four node types (\textsc{Endpoint}, \textsc{Middleware}, \textsc{Service}, \textsc{Data\_Access}) connected by three edge types (\textsc{Call}, \textsc{Identifier\_Flow}, \textsc{Tenant\_Flow}), and encodes a 36-dimensional feature vector per node where each dimension corresponds directly to a specific, interpretable authorization property.

This design yields two immediate advantages. First, because every feature is interpretable and every node represents a concrete architectural component, the PFG provides \emph{inherent interpretability}: a developer can inspect the graph and immediately understand the authorization posture of any endpoint without post-hoc explanation tools. Second, because the graph is compact---typically 5 to 50~nodes per application, compared to thousands for a CPG---the GNN operates on a focused, noise-free representation where every element is security-relevant.

We complement the GNN classifier with a \textbf{structural validator}: a set of five formally defined detection rules that trace explicit paths from endpoints to data-access nodes, producing human-readable explanations. Unlike GNNExplainer, which learns attribution masks post-hoc, our validator provides explanations \emph{by construction}---the affected endpoint, the vulnerability subtype, and the missing check along the path are derived directly from the graph structure.

\medskip

The remainder of this chapter is organized as follows. Section~\ref{sec:problem} formalizes the problem: the threat model, the taxonomy of five targeted vulnerability subtypes, and three research questions. Section~\ref{sec:positioning} positions our approach against NLP-based and CPG-based alternatives, articulating the design rationale. Section~\ref{sec:pfg} defines the PFG formally. Section~\ref{sec:gnn} presents the GNN architecture. Section~\ref{sec:validator} describes the structural validator and its complementarity with the learned model. Finally, Section~\ref{sec:ch3-conclusion} summarizes the design choices.


% ============================================================
\section{Problem Formalization}
\label{sec:problem}

\subsection{Threat Model and Scope}
\label{sec:threat-model}

We define the threat model for our framework as follows. The \textbf{attacker} is a user who possesses valid credentials---that is, they can authenticate to the system---but attempts to access resources or perform operations beyond their authorized privilege level. This corresponds to the insider threat or compromised-account scenario, which is the most common vector for authorization vulnerabilities in production web applications.

The \textbf{assets under protection} are REST API endpoints that expose access to persistent data via database queries. Each endpoint may handle one or more HTTP methods (GET, POST, PUT, DELETE, PATCH) and may accept route parameters that identify specific objects (e.g., \texttt{/api/users/\{pk\}/}).

We define three \textbf{trust boundaries} through which every legitimate request must pass:

\begin{enumerate}
    \item \textbf{Authentication boundary:} verifies that the requester possesses a valid identity (session token, JWT, API key).
    \item \textbf{Authorization boundary:} verifies that the authenticated identity has sufficient privilege for the requested operation (role check, permission class, access control list).
    \item \textbf{Object-level boundary:} verifies that the authenticated, authorized user is permitted to access the \emph{specific} object identified by the request parameters (ownership check, tenant scoping, object-level permission).
\end{enumerate}

An authorization vulnerability exists when a request can traverse from an endpoint to a data-access operation while bypassing one or more of these boundaries.

\textbf{Scope.} Our framework targets server-side web applications built with Django/Django REST Framework (Python) and Spring Boot (Java) that expose REST APIs. We explicitly exclude: client-side vulnerabilities, injection flaws (SQL injection, XSS, CSRF), cryptographic weaknesses, network-level attacks, and denial-of-service. These exclusions are deliberate: our graph representation is designed specifically for the structural properties of authorization, and extending it to other vulnerability classes would require fundamentally different node types and features.

\subsection{Taxonomy of Authorization Vulnerabilities}
\label{sec:taxonomy}

We target five distinct subtypes of authorization vulnerabilities, each corresponding to the bypass of a specific trust boundary. Table~\ref{tab:taxonomy} presents the formal taxonomy.

\begin{table}[htbp]
\centering
\caption{Taxonomy of targeted authorization vulnerability subtypes.}
\label{tab:taxonomy}
\begin{tabular}{@{}p{4.2cm}p{1.5cm}p{7.8cm}@{}}
\toprule
\textbf{Subtype} & \textbf{CWE} & \textbf{Definition} \\
\midrule
Missing Authentication & CWE-306 & An endpoint is reachable without any identity verification before performing a data-access operation. \\
\addlinespace
Missing Role Check & CWE-862 & An endpoint verifies the requester's identity but does not enforce a privilege-level check before performing a sensitive operation (e.g., administrative action, write to shared resource). \\
\addlinespace
Insecure Direct Object Reference (IDOR) & CWE-639 & A route parameter identifying a specific object flows to a data-access operation without verification that the requester owns or is authorized to access that object. \\
\addlinespace
Tenant Isolation Failure & CWE-284 & In a multi-tenant application, data belonging to one tenant is reachable by a user of another tenant due to missing tenant-scoping predicates in queries. \\
\addlinespace
Object-Level Permission Bypass & CWE-285 & The application defines object-level permissions (e.g., Django's \texttt{has\_object\_permission}), but the access path does not enforce them before returning the object. \\
\bottomrule
\end{tabular}
\end{table}

In addition to these five vulnerability subtypes, our classification model recognizes three supplementary labels: \textsc{Safe} (no vulnerability detected), \textsc{Combined} (multiple subtypes apply simultaneously to the same endpoint), and \textsc{Intentionally\_Public} (the endpoint is explicitly marked as public, e.g., via \texttt{AllowAny} or \texttt{@PermitAll}).

Each subtype maps to the OWASP API Security Top~10 (2023 edition): Missing Authentication and Missing Role Check correspond to \textbf{API1:2023---Broken Object Level Authorization} and \textbf{API5:2023---Broken Function Level Authorization}; IDOR maps to \textbf{API1:2023}; Tenant Isolation Failure relates to \textbf{API1:2023} in multi-tenant contexts; and Object-Level Permission Bypass corresponds to \textbf{API1:2023} when object permissions exist but are not enforced.

\textbf{Justification of scope.} Standard vulnerable-application benchmarks such as OWASP~WebGoat and the Damn Vulnerable Web Application~(DVWA) are unsuitable for our work. These benchmarks target \emph{injection} vulnerabilities---SQL injection, cross-site scripting, cross-site request forgery---presented as intentionally obvious single-file exercises. Authorization vulnerabilities differ fundamentally: they are \emph{structural}, manifesting as the \emph{absence} of checks along execution paths rather than the \emph{presence} of exploitable input handling. Moreover, these benchmarks are monolithic educational applications, not production code with hundreds of endpoints, ViewSets, permission classes, and layered middleware stacks. Our evaluation therefore relies on real-world CVE-backed vulnerabilities from production open-source applications (detailed in Chapter~4).

\subsection{Research Questions}
\label{sec:research-questions}

Building on the gaps identified in Part~I, we formulate three research questions:

\begin{description}
    \item[RQ1 (Detection Effectiveness):] Can a domain-specific graph representation of authorization flows---the Program Flow Graph---effectively detect access-control vulnerabilities in real-world web applications? \\
    \emph{Metric:} Precision, Recall, and F1-score on a benchmark of 75~real-world CVEs from production Django and Spring applications.

    \item[RQ2 (Inherent Interpretability):] Does the PFG framework provide interpretable detection results without requiring post-hoc explanation tools such as GNNExplainer? \\
    \emph{Metric:} Each reported finding includes the affected endpoint identifier, the vulnerability subtype, the data-access target, and a path-based natural-language explanation.

    \item[RQ3 (Comparison with Existing Tools):] How does GNN-based detection on the PFG compare to rule-based static analysis and industry-standard SAST tools (Semgrep, CodeQL) on the same benchmark? \\
    \emph{Metric:} Side-by-side Precision, Recall, and F1-score using identical ground truth.
\end{description}


% ============================================================
\section{Positioning: Why Not Existing Approaches}
\label{sec:positioning}

Before presenting our solution, we examine why the two dominant paradigms---NLP-based and CPG-based vulnerability detection---are structurally ill-suited for authorization vulnerabilities, and articulate the design principles that guided our approach.

\subsection{Limitations of NLP-Based Approaches}
\label{sec:nlp-limits}

Pre-trained code language models such as CodeBERT~\cite{feng2020codebert} and VulBERTa~\cite{hanif2022vulberta} encode source code into dense vector representations by processing token sequences through Transformer architectures. The resulting embeddings (typically 768~dimensions) capture lexical and statistical patterns in code, enabling downstream classifiers to identify functions that ``look like'' previously-seen vulnerable code.

For authorization vulnerabilities, this paradigm has three fundamental limitations:

\begin{enumerate}
    \item \textbf{Structural blindness.} Authorization is determined by the \emph{presence or absence} of checks along the request-processing chain---a structural property spanning multiple functions and modules. A token-level model processes a single function at a time and cannot reason about whether an authentication middleware exists upstream, whether a permission class is inherited from a base class, or whether a queryset filter enforces ownership downstream.

    \item \textbf{Opacity of embeddings.} The 768-dimensional CodeBERT embedding is opaque: one cannot inspect a specific dimension and determine whether it encodes ``this function has a \texttt{@login\_required} decorator.'' When the model flags a function, it cannot explain \emph{why}---it can only report a confidence score. For security engineers who must triage findings, this lack of explanation is a serious practical limitation.

    \item \textbf{Mismatch with the vulnerability signal.} Authorization vulnerabilities are about what is \emph{missing}, not what is \emph{present}. Two functions may be lexically identical except that one resides behind an authentication middleware and the other does not. A token-level model, operating on the function's source alone, cannot distinguish these cases.
\end{enumerate}

\subsection{Limitations of CPG-Based Approaches with GNNExplainer}
\label{sec:cpg-limits}

Code Property Graphs (CPGs) address structural blindness by representing code as a graph that merges AST, CFG, and PDG edges~\cite{yamaguchi2014modeling}. Tools such as Joern~\cite{yamaguchi2014modeling} construct CPGs automatically, and GNN models trained on CPG subgraphs have shown promise for vulnerability detection~\cite{zhou2019devign, li2021vulnerability}.

However, for authorization vulnerability detection, CPGs present distinct challenges:

\begin{enumerate}
    \item \textbf{Excessive granularity.} A typical function generates hundreds to thousands of CPG nodes (one per AST node, plus CFG and PDG edges). For a function with 50~lines of code, the CPG may contain 500+ nodes---most representing assignments, literals, operators, and control-flow joins that are irrelevant to the authorization question. The relevant information (``is there a permission check?'') is buried in this noise.

    \item \textbf{Post-hoc interpretability.} GNNExplainer~\cite{ying2019gnnexplainer} was developed precisely because CPG-based predictions are uninterpretable by default. It learns an edge mask that maximizes mutual information with the prediction, identifying a ``critical subgraph.'' However, this subgraph still requires expert interpretation: knowing that nodes~42, 87, and 156 were important does not directly tell a developer which authorization check is missing or where to add it.

    \item \textbf{Irrelevant abstraction level.} Authorization decisions in web frameworks operate at the \emph{architectural} level: middleware chains, decorator stacks, class hierarchies, queryset methods. The CPG, by contrast, operates at the \emph{syntactic} level: individual statements, expressions, and data-flow edges. There is a fundamental mismatch between the abstraction level at which the vulnerability exists and the abstraction level at which the CPG represents code.
\end{enumerate}

\subsection{Our Design Rationale}
\label{sec:rationale}

The limitations above motivate three design principles for our framework:

\begin{enumerate}
    \item \textbf{Model only what matters.} Rather than representing all code structure (CPG) or treating code as text (CodeBERT), we represent exclusively the authorization-relevant architecture: the chain from HTTP entry point, through any guards, to data access. This yields a graph of 5--50 nodes where every element is semantically meaningful.

    \item \textbf{Encode security properties directly.} Rather than learning implicit representations from code tokens, we extract explicit, interpretable features: \texttt{requires\_auth}, \texttt{has\_role\_check}, \texttt{has\_ownership\_predicate}, \texttt{has\_tenant\_scope}. Each of the 36~features in our vector can be inspected by a human and directly corresponds to an authorization-relevant property.

    \item \textbf{Build explanation into the structure.} Rather than applying post-hoc explainability tools, we design the graph so that its structure \emph{is} the explanation. The path from an endpoint node to a data-access node, with or without intervening middleware, directly encodes whether and how authorization is enforced. A structural validator can trace this path and produce a human-readable explanation without any additional learned component.
\end{enumerate}

The role of the GNN in our architecture is then well-defined: it learns \emph{interaction effects} and \emph{context-dependent patterns} that cannot be captured by explicit rules alone. For example, the combination of a missing role check, the presence of a write operation, and the absence of ownership enforcement may collectively indicate a vulnerability that no single rule would flag. The GNN provides the learning \emph{ceiling}; the rules provide the precision \emph{floor}.


% ============================================================
\section{The Program Flow Graph}
\label{sec:pfg}

\subsection{Formal Definition}
\label{sec:pfg-definition}

We define the Program Flow Graph as a directed, typed, attributed graph:

\begin{equation}
\label{eq:pfg}
G = (V, E, \tau_V, \tau_E, \varphi_V, \varphi_E)
\end{equation}

\noindent where:
\begin{itemize}
    \item $V$ is the set of nodes representing architectural components;
    \item $E \subseteq V \times V$ is the set of directed edges representing relationships between components;
    \item $\tau_V: V \to \{\textsc{Endpoint}, \textsc{Middleware}, \textsc{Service}, \textsc{Data\_Access}\}$ assigns a type to each node;
    \item $\tau_E: E \to \{\textsc{Call}, \textsc{Identifier\_Flow}, \textsc{Tenant\_Flow}\}$ assigns a type to each edge;
    \item $\varphi_V: V \to \mathbb{R}^{36}$ maps each node to a 36-dimensional feature vector;
    \item $\varphi_E: E \to \mathbb{R}^{3}$ maps each edge to a 3-dimensional feature vector (one-hot encoding of edge type).
\end{itemize}

The PFG has the following key properties:

\begin{itemize}
    \item \textbf{Directed:} edges follow the direction of request processing, from endpoint toward data access.
    \item \textbf{Compact:} a typical application yields 5--50 nodes (compared to thousands for a CPG of the same application).
    \item \textbf{Semantically grounded:} every node corresponds to a concrete, identifiable component in the source code (a specific URL handler, a specific ORM query, a specific middleware).
    \item \textbf{Schema-validated:} all PFG instances conform to a JSON schema (version~1.0) that enforces structural constraints.
\end{itemize}

\subsection{Node Types and Attributes}
\label{sec:node-types}

The PFG defines four node types, each modeling a distinct layer in the authorization architecture of a web application.

\subsubsection{Endpoint Node}

An \textsc{Endpoint} node represents an HTTP route handler---the entry point for client requests. It carries the following attributes:

\begin{itemize}
    \item \texttt{http\_method} $\in$ \{GET, POST, PUT, DELETE, PATCH\}: the HTTP method handled.
    \item \texttt{route}: the URL pattern (e.g., \texttt{/api/users/\{pk\}/}).
    \item \texttt{route\_params}: list of path parameters extracted from the route.
    \item \texttt{requires\_auth}: whether the endpoint enforces authentication.
    \item \texttt{required\_role} $\in$ \{None, User, Moderator, Admin, Superadmin\}: minimum privilege level required.
    \item \texttt{ownership\_check\_present}: whether an ownership verification exists on any path to data access.
    \item \texttt{has\_route\_params}: whether the route contains object-identifying parameters.
    \item \texttt{reachable\_data\_count}: number of data-access nodes reachable from this endpoint.
    \item \texttt{has\_id\_flow}: whether a route parameter flows to a data-access argument.
    \item \texttt{unprotected\_data\_reachable}: whether any data access is reachable without passing through a guard.
    \item Authorization signal attributes: \texttt{role\_check\_present}, \texttt{model\_permission\_check\_present}, \texttt{object\_permission\_check\_present}, \texttt{tenant\_check\_present}, \texttt{has\_tenant\_scope}, \texttt{unscoped\_tenant\_data\_reachable}, \texttt{object\_permission\_data\_reachable}.
\end{itemize}

\subsubsection{Middleware Node}

A \textsc{Middleware} node represents an authorization guard---any component that enforces a security check. It carries:

\begin{itemize}
    \item \texttt{guard\_type} $\in$ \{\textsc{Authentication}, \textsc{Role\_Check}, \textsc{Ownership\_Check}, \textsc{Tenant\_Check}, \textsc{Object\_Permission\_Check}\}: the category of authorization enforced.
    \item \texttt{guard\_position}: the ordinal position in the middleware chain (lower values execute first).
\end{itemize}

In the PFG, middleware nodes appear on paths between endpoints and data-access nodes. Their presence or absence on a given path directly determines the authorization posture of the corresponding endpoint.

\subsubsection{Service Node}

A \textsc{Service} node represents a business-logic component---a function, method, or class that performs computation between the endpoint and the data layer. It carries:

\begin{itemize}
    \item \texttt{name}: the fully qualified name (qualname) of the function or class.
\end{itemize}

Service nodes enable the PFG to model indirect data access: an endpoint may not query the database directly but instead delegate to a service function that performs the actual ORM call.

\subsubsection{Data Access Node}

A \textsc{Data\_Access} node represents a database operation. It carries:

\begin{itemize}
    \item \texttt{operation} $\in$ \{Read, Write, Update, Delete\}: the type of database operation.
    \item \texttt{target\_entity}: the database model or table accessed.
    \item \texttt{has\_ownership\_predicate}: whether the query filters by the requesting user (e.g., \texttt{.filter(owner=request.user)}).
    \item \texttt{has\_tenant\_predicate}: whether the query enforces tenant scoping.
    \item \texttt{has\_object\_permission\_check}: whether object-level permission is verified before returning the result.
\end{itemize}

\subsection{Edge Types and Semantics}
\label{sec:edge-types}

The PFG defines three edge types, each capturing a distinct relationship between components:

\begin{description}
    \item[\textsc{Call}] represents a function invocation relationship. An edge $(u, v)$ of type \textsc{Call} indicates that component~$u$ invokes component~$v$ during request processing. This models the call chain from endpoint through services to data access, as well as the invocation of middleware guards.

    \item[\textsc{Identifier\_Flow}] represents the propagation of a user-supplied identifier (typically a route parameter such as a primary key) from the endpoint to a data-access argument. This edge type is critical for IDOR detection: if an identifier flows from the route to a database query without passing through an ownership check, the endpoint may be vulnerable.

    \item[\textsc{Tenant\_Flow}] represents the propagation of tenant-scoping information (organization ID, workspace identifier) through the request-processing chain. This edge type enables detection of tenant isolation failures by tracking whether tenant context reaches the data-access layer.
\end{description}

\subsection{Feature Engineering}
\label{sec:features}

Each node in the PFG is encoded as a 36-dimensional real-valued feature vector. The vector is composed of the following segments:

\begin{table}[htbp]
\centering
\caption{PFG node feature vector layout (36 dimensions).}
\label{tab:features}
\begin{tabular}{@{}cll@{}}
\toprule
\textbf{Indices} & \textbf{Feature} & \textbf{Encoding} \\
\midrule
$[0{:}4]$ & Node type & One-hot (4 classes) \\
$[4{:}9]$ & HTTP method & One-hot (5 methods) \\
$[9]$ & Requires authentication & Binary \\
$[10{:}15]$ & Required role & One-hot (5 levels) \\
$[15]$ & Ownership check present & Binary \\
$[16{:}19]$ & Guard type & Partial one-hot (3 dims) \\
$[19{:}23]$ & Operation type & One-hot (4 operations) \\
$[23]$ & Has ownership predicate & Binary \\
$[24]$ & Has route parameters & Binary \\
$[25]$ & Reachable data count & Normalized float \\
$[26]$ & Has identifier flow & Binary \\
$[27]$ & Unprotected data reachable & Binary \\
$[28]$ & Role check present & Binary \\
$[29]$ & Model permission check present & Binary \\
$[30]$ & Object permission check present & Binary \\
$[31]$ & Tenant check present & Binary \\
$[32]$ & Has tenant scope & Binary \\
$[33]$ & Unscoped tenant data reachable & Binary \\
$[34]$ & Object permission data reachable & Binary \\
$[35]$ & Guard signal (composite) & Float $\in \{0, 0.33, 0.66, 1.0\}$ \\
\bottomrule
\end{tabular}
\end{table}

Each edge is encoded as a 3-dimensional one-hot vector corresponding to its type: $[1,0,0]$ for \textsc{Call}, $[0,1,0]$ for \textsc{Identifier\_Flow}, and $[0,0,1]$ for \textsc{Tenant\_Flow}.

\textbf{Design argument.} Every dimension in this feature vector has a clear, human-interpretable meaning tied to a specific authorization property. A security engineer can inspect feature index~9 and immediately determine whether the endpoint requires authentication; they can inspect index~23 and determine whether the data-access query filters by ownership. This stands in direct contrast to CodeBERT embeddings, where no individual dimension carries interpretable semantic content. Our features require no GPU computation, no pre-trained language model, and no inference pipeline---they are computed directly by static analysis of the source code.

\subsection{PFG Construction Pipeline}
\label{sec:construction}

The PFG is constructed through a seven-stage pipeline, executed by framework-specific extractors (detailed in Chapter~4):

\begin{enumerate}
    \item \textbf{Endpoint discovery:} identify all HTTP route handlers by analyzing URL configuration files, router registrations, and annotated controller classes.
    \item \textbf{Authentication/authorization annotation:} for each endpoint, determine whether authentication is required, which role (if any) is enforced, and whether the endpoint is intentionally public. This is achieved by inspecting decorators, class-level attributes, mixin inheritance, and framework settings.
    \item \textbf{Call graph construction:} perform intra-project AST traversal with import alias resolution to identify which functions each endpoint invokes, transitively up to a bounded depth.
    \item \textbf{Data access detection:} identify all database operations (ORM calls, repository method invocations) within reachable functions, classifying each by operation type and ownership enforcement.
    \item \textbf{Identifier flow tracking:} perform intra-procedural taint propagation from route parameters through assignments and function arguments to identify which data-access operations receive user-supplied identifiers.
    \item \textbf{Graph assembly:} create nodes for each discovered component, insert edges according to the call graph and flow analyses, compute derived attributes, and insert \textsc{Middleware} nodes for detected guards.
    \item \textbf{Schema validation:} validate the resulting PFG against the JSON schema to ensure structural integrity before downstream consumption.
\end{enumerate}


% ============================================================
\section{GNN-Based Vulnerability Detection}
\label{sec:gnn}

\subsection{Motivation: Why a GNN on Top of Rules?}
\label{sec:gnn-motivation}

The structural validator (Section~\ref{sec:validator}) detects vulnerabilities through explicit pattern matching: if specific conditions hold on the PFG (e.g., no authentication guard on a path to data access), a finding is emitted. This approach provides high precision on well-defined patterns but is inherently limited to what can be formally specified.

Real-world authorization, however, is \emph{contextual}. An endpoint without \texttt{@login\_required} may nonetheless be protected by a class-level mixin inherited from a parent view. A missing ownership check may be acceptable if the endpoint only reads public, non-sensitive data. The \emph{combination} of multiple weak signals---a write operation targeting a shared resource, the absence of a role check, and the presence of an identifier flow---may collectively indicate a vulnerability that no single rule would flag.

A GNN addresses these limitations by learning from labeled data:
\begin{itemize}
    \item It captures \textbf{interaction effects} between features and graph topology that rules cannot express.
    \item It handles \textbf{context-dependent decisions} where the same local pattern is vulnerable in one architectural context but safe in another.
    \item It provides a learning \textbf{ceiling} above what rules achieve, while rules provide a precision \textbf{floor} that the GNN can never fall below (in the combined architecture).
\end{itemize}

\subsection{PathAnalyzer: Primary Multi-Class Model}
\label{sec:path-analyzer}

Our primary detection model, \textbf{PathAnalyzer}, is a graph-level classifier that takes an entire PFG as input and classifies a specific endpoint within it. The architecture comprises four stages:

\subsubsection{Input Projection}

The 36-dimensional node feature vector is projected into the model's hidden space:
\begin{equation}
\mathbf{h}_i^{(0)} = \text{Linear}_{36 \to d}(\varphi_V(v_i))
\end{equation}
where $d = 64$ is the hidden dimension.

\subsubsection{Message Passing (GATv2Conv)}

We employ three layers of GATv2~\cite{brody2021attentive} convolution with multi-head attention:
\begin{equation}
\mathbf{h}_i^{(\ell+1)} = \Big\Vert_{k=1}^{K} \sum_{j \in \mathcal{N}(i)} \alpha_{ij}^{k} \mathbf{W}^k \mathbf{h}_j^{(\ell)}
\end{equation}
where $K=4$ attention heads, each producing $d/K = 16$-dimensional outputs that are concatenated to yield a $d$-dimensional representation. Edge features (3-dim) are integrated into the attention computation, allowing the model to attend differently to \textsc{Call}, \textsc{Identifier\_Flow}, and \textsc{Tenant\_Flow} edges.

\textbf{Justification of GATv2 over GCN/GraphSAGE:} The attention mechanism captures \emph{asymmetric} relationships---an authentication middleware protects all downstream nodes but not upstream ones. Multi-head attention captures multiple aspects of authorization relationships simultaneously (e.g., one head may attend to guards, another to data-access patterns).

\subsubsection{Target-Aware Attention Pooling}

Unlike standard graph classification that applies global mean or max pooling, we employ a \emph{target-aware} attention mechanism. Given a target endpoint node~$t$ whose vulnerability status we wish to classify:

\begin{align}
\mathbf{s}_i &= \text{Linear}_{2d \to d}\left([\mathbf{h}_i^{(L)} \Vert \mathbf{h}_t^{(L)}]\right) \\
a_i &= \frac{\exp(\text{tanh}(\text{Linear}_{d \to 1}(\mathbf{s}_i)))}{\sum_{j} \exp(\text{tanh}(\text{Linear}_{d \to 1}(\mathbf{s}_j)))} \\
\mathbf{c} &= \sum_{i \in V} a_i \cdot \mathbf{h}_i^{(L)}
\end{align}

This mechanism computes attention weights relative to the target endpoint, enabling the model to focus on the specific context relevant to classifying \emph{that} endpoint. Different endpoints in the same graph may have different vulnerability statuses; target-aware pooling ensures the representation is conditioned on the endpoint under examination.

\subsubsection{Classification Head}

The final representation concatenates the target endpoint embedding with the context vector:
\begin{equation}
\hat{y} = \text{Linear}_{d \to C}\left(\text{ReLU}\left(\text{Linear}_{2d \to d}\left([\mathbf{h}_t^{(L)} \Vert \mathbf{c}]\right)\right)\right)
\end{equation}
where $C = 8$ is the number of classes: \textsc{Safe}, \textsc{Missing\_Auth}, \textsc{Missing\_Role}, \textsc{IDOR}, \textsc{Tenant\_Isolation\_Failure}, \textsc{Object\_Level\_Permission\_Bypass}, \textsc{Combined}, and \textsc{Intentionally\_Public}.

\subsection{PfgAnomalyDetector: Binary Screening Model}
\label{sec:anomaly-detector}

For scenarios requiring fast initial screening, we provide a binary classifier, \textbf{PfgAnomalyDetector}, based on GINEConv~\cite{hu2019strategies}:

\begin{itemize}
    \item Three GINEConv layers with inner MLPs (Linear $\to$ ReLU $\to$ Linear), edge dimension~3.
    \item Global pooling: concatenation of mean and max pooling over all nodes, yielding a $2d$-dimensional graph-level representation.
    \item Classification head: Linear$(2d \to d)$ $\to$ ReLU $\to$ Dropout $\to$ Linear$(d \to 2)$.
    \item Output: binary logits (vulnerable vs.\ benign).
\end{itemize}

This model serves as a lightweight first-pass detector: it identifies which endpoints warrant detailed multi-class analysis by PathAnalyzer.

\subsection{Training Objective and Class Imbalance}
\label{sec:training-objective}

Both models are trained with class-weighted cross-entropy loss:
\begin{equation}
\mathcal{L} = -\sum_{i=1}^{N} w_{y_i} \log \frac{\exp(\hat{y}_i[y_i])}{\sum_{c=1}^{C} \exp(\hat{y}_i[c])}
\end{equation}
where $w_c$ is the inverse frequency weight for class~$c$. This addresses the severe class imbalance inherent in authorization vulnerability detection: the vast majority of endpoints are safe, and vulnerable endpoints are rare.

During inference, per-subtype confidence thresholds (rather than a uniform 0.5 threshold) are applied, calibrated on the validation set to optimize the precision-recall trade-off for each vulnerability subtype independently.


% ============================================================
\section{Structural Validator}
\label{sec:validator}

\subsection{Five Detection Rules}
\label{sec:rules}

The structural validator implements five formally defined detection rules, each corresponding to one vulnerability subtype. A rule fires when specific structural conditions hold on the PFG.

\begin{description}
    \item[R1 --- Missing Authentication.] An endpoint node with \texttt{requires\_auth} $=$ \textit{false} has a reachable path to at least one \textsc{Data\_Access} node, and the endpoint is not marked as intentionally public.
    \begin{equation}
    \text{R1}: \neg \text{auth}(e) \wedge \exists\, p \in \text{paths}(e, d) \wedge \neg \text{public}(e) \implies \text{MISSING\_AUTH}
    \end{equation}

    \item[R2 --- Missing Role Check.] An endpoint reaches a \textsc{Data\_Access} node that performs a privileged operation (administrative action, or write/delete on a shared resource), and no \textsc{Middleware} node of type \textsc{Role\_Check} exists on any path from the endpoint to that data access.
    \begin{equation}
    \text{R2}: \text{priv}(d) \wedge \nexists\, m \in p : \tau_V(m) = \textsc{Mid} \wedge \text{guard}(m) = \textsc{Role} \implies \text{MISSING\_ROLE}
    \end{equation}

    \item[R3 --- IDOR.] The endpoint has route parameters, an \textsc{Identifier\_Flow} edge connects a parameter to a \textsc{Data\_Access} node operating on an ownership-bearing entity, and no ownership check exists on any path (neither a \textsc{Middleware} node of type \textsc{Ownership\_Check} nor a data-access predicate filtering by the requesting user).
    \begin{equation}
    \text{R3}: \text{params}(e) \wedge \text{id\_flow}(e, d) \wedge \neg \text{own\_check}(p) \wedge \neg \text{own\_pred}(d) \implies \text{IDOR}
    \end{equation}

    \item[R4 --- Tenant Isolation Failure.] Tenant-scoped data is reachable, but neither the endpoint nor the data-access node enforces tenant scoping.
    \begin{equation}
    \text{R4}: \text{tenant\_data}(d) \wedge \neg \text{tenant\_scope}(e) \wedge \neg \text{tenant\_pred}(d) \implies \text{TENANT\_ISO}
    \end{equation}

    \item[R5 --- Object-Level Permission Bypass.] Object-level permissions are defined for the accessed model, the access uses a route parameter, but no object-permission check is enforced on the path.
    \begin{equation}
    \text{R5}: \text{obj\_perm}(d) \wedge \text{params}(e) \wedge \neg \text{obj\_check}(p) \implies \text{OBJ\_PERM\_BYPASS}
    \end{equation}
\end{description}

\subsection{Role as Interpretability Layer}
\label{sec:interpretability}

Each rule, when triggered, produces a \textbf{Finding} comprising four elements:

\begin{enumerate}
    \item \texttt{endpoint\_id}: the fully qualified name of the affected endpoint.
    \item \texttt{subtype}: the vulnerability category from our taxonomy.
    \item \texttt{data\_access\_id}: the specific database operation that is unprotected.
    \item \texttt{explanation}: a human-readable natural-language description of the violation.
\end{enumerate}

For example, a typical finding reads: \emph{``Endpoint \texttt{UserViewSet.retrieve} reaches \texttt{User.objects.get(pk=pk)} via identifier\_flow from route parameter \texttt{pk} without ownership verification on the path.''} This explanation directly answers the three questions a developer needs:

\begin{itemize}
    \item \textbf{Where:} which endpoint and route is affected.
    \item \textbf{What:} which vulnerability subtype applies.
    \item \textbf{Why:} which check is missing and which data access is exposed.
\end{itemize}

This constitutes our answer to \textbf{RQ2}: the PFG provides inherent interpretability because its structure \emph{is} the explanation. The structural validator merely reads and verbalizes the graph topology---no learned masks, no gradient attribution, no post-hoc approximation.

\subsection{Complementarity of Rules and GNN}
\label{sec:complementarity}

The structural validator and the GNN are complementary components in our detection architecture:

\begin{itemize}
    \item \textbf{Rules} provide high precision on explicitly formalizable patterns. They never produce false positives on patterns they are designed to detect (given a correct PFG). However, they cannot detect vulnerabilities that do not match their explicitly coded conditions.

    \item \textbf{The GNN} learns from data and captures implicit patterns, interaction effects, and context-dependent relationships that resist formal specification. It can generalize to vulnerability patterns not seen during rule design. However, it may produce false positives on unfamiliar graph topologies.

    \item \textbf{During training,} the rules serve a dual role: they provide weak supervision labels for unlabeled endpoints (generating pseudo-labeled training samples), expanding the training set beyond the CVE-labeled ground truth.

    \item \textbf{During inference,} the architecture supports multiple combination strategies: the GNN proposes candidates that the validator confirms and explains (high-precision mode), or both operate independently and their findings are unioned (high-recall mode).
\end{itemize}

The ablation study in Chapter~5 quantifies the contribution of each component independently and in combination.


% ============================================================
\section{Conclusion}
\label{sec:ch3-conclusion}

This chapter presented the design of the PFG framework for authorization vulnerability detection, founded on five deliberate architectural choices:

\begin{enumerate}
    \item \textbf{Domain-specific graph representation:} the PFG models only authorization-relevant architecture (4~node types, 3~edge types), yielding compact, semantically grounded graphs that provide inherent interpretability without post-hoc tools.

    \item \textbf{Interpretable feature engineering:} the 36-dimensional node feature vector encodes authorization properties directly, eliminating the need for opaque learned embeddings (CodeBERT) and enabling human inspection of every dimension.

    \item \textbf{GATv2 with target-aware attention pooling:} the PathAnalyzer model leverages multi-head attention to capture asymmetric authorization relationships and conditions its prediction on the specific endpoint under examination.

    \item \textbf{Structural validator:} five formally defined rules provide human-readable explanations by construction, answering where, what, and why for each finding.

    \item \textbf{Rules--GNN complementarity:} explicit rules provide a precision floor and interpretability, while the GNN raises the detection ceiling by learning patterns beyond formal specification.
\end{enumerate}

The next chapter details the concrete implementation of this framework: the framework-specific extractors for Django and Spring Boot, the construction of a real-world evaluation dataset from 75~CVE entries, the training pipeline with multiple data strategies, and the demonstration system.


% ============================================================
% ============================================================
% CHAPTER 4 — IMPLEMENTATION & DATASET
% ============================================================
% ============================================================

\chapter{Implementation \& Dataset}
\label{chap:implementation}

\section{Introduction}
\label{sec:impl-intro}

Chapter~3 defined the PFG framework at the \emph{design} level: the graph schema, the GNN architecture, and the structural validator. This chapter bridges the gap between design and practice by presenting the concrete implementation. We address four questions: (1)~how are PFGs extracted from real Django/DRF and Spring Boot codebases? (2)~how was the evaluation dataset constructed from real-world CVE entries? (3)~how is the training pipeline organized to handle data scarcity and class imbalance? and (4)~how does the demonstration system expose the framework to end users?

The chapter is organized as follows. Section~\ref{sec:tech-stack} presents the technology stack and justifies each tool choice. Section~\ref{sec:extractors} details the framework-specific extractors. Section~\ref{sec:dataset} describes the dataset construction process. Section~\ref{sec:training-pipeline} explains the training pipeline and its three configurations. Section~\ref{sec:demo} presents the demonstration system. Section~\ref{sec:ch4-conclusion} concludes.


% ============================================================
\section{Technology Stack}
\label{sec:tech-stack}

Table~\ref{tab:tech-stack} summarizes the technologies employed in the implementation and the rationale for each choice.

\begin{table}[htbp]
\centering
\caption{Technology stack and justifications.}
\label{tab:tech-stack}
\begin{tabular}{@{}p{3.2cm}p{3.5cm}p{7cm}@{}}
\toprule
\textbf{Component} & \textbf{Technology} & \textbf{Justification} \\
\midrule
Django code parsing & Python \texttt{ast} module & Native standard-library module; provides full Abstract Syntax Tree access without external dependencies. \\
\addlinespace
Spring code parsing & Regex-based extractor & Pragmatic choice avoiding a full Java parser dependency; sufficient for annotation-driven Spring patterns. \\
\addlinespace
Graph representation & Pydantic models + JSON & Typed data models with runtime validation; human-readable serialization; portable across pipeline stages. \\
\addlinespace
Schema validation & JSON Schema (v1.0) & Ensures structural integrity of PFGs before downstream consumption; catches extractor bugs early. \\
\addlinespace
Graph ML framework & PyTorch Geometric (PyG) & Native support for GATv2Conv, GINEConv, heterogeneous batching, and edge-feature integration. \\
\addlinespace
Optimization & AdamW optimizer & Decoupled weight decay; standard for GNN training~[citation needed]. \\
\addlinespace
Backend API & FastAPI + WebSocket & Asynchronous request handling; native WebSocket support for real-time streaming of analysis progress. \\
\addlinespace
Frontend & Next.js~14 (React) & Server-side rendering; component-based architecture; modern developer experience. \\
\addlinespace
Graph visualization & D3.js (frontend); Graphviz DOT (reports) & Interactive force-directed layouts for the demo; static publication-quality graphs for reports. \\
\bottomrule
\end{tabular}
\end{table}

\textbf{Why not Joern?} Joern~\cite{yamaguchi2014modeling} is an industrial Code Property Graph builder that constructs generic CPGs (AST + CFG + PDG). We chose not to use it for three reasons. First, Joern produces graphs with thousands of nodes per function, whereas our PFG deliberately contains only authorization-relevant components (5--50 nodes per application). Second, Joern cannot compute the authorization-specific features we require (\texttt{requires\_auth}, \texttt{has\_ownership\_predicate}, \texttt{permission\_classes} inheritance) because these depend on framework-specific semantic analysis. Third, Joern's JVM runtime introduces a heavyweight dependency unsuitable for our lightweight, Python-native pipeline.

\textbf{Why not CodeBERT for features?} As argued in Chapter~3 (Section~\ref{sec:nlp-limits}), CodeBERT embeddings are opaque 768-dimensional vectors that encode generic code semantics. Our 36-dimensional feature vector is computed by direct static analysis, requires no GPU, and is fully interpretable. Each dimension answers a specific authorization question.


% ============================================================
\section{Framework-Specific Extractors}
\label{sec:extractors}

The PFG construction pipeline (Section~\ref{sec:construction}) is implemented by two framework-specific extractors: one for Django/Django REST Framework (Python) and one for Spring Boot (Java). Each extractor implements the seven-stage pipeline defined in Chapter~3 using techniques appropriate to the target framework's conventions and language.

\subsection{Django Extractor}
\label{sec:django-extractor}

The Django extractor is an AST-based static analyzer implemented in Python. It comprises six modules, each responsible for one aspect of PFG construction.

\subsubsection{Endpoint Discovery}

The endpoint discovery module identifies all HTTP route handlers by analyzing Django's URL configuration system. It processes:

\begin{itemize}
    \item \textbf{Standard URL patterns:} \texttt{path()}, \texttt{re\_path()}, and nested \texttt{include()} directives in \texttt{urls.py} files.
    \item \textbf{DRF router registrations:} \texttt{DefaultRouter.register()} calls that automatically generate CRUD endpoints for ViewSets.
    \item \textbf{Class-Based View (CBV) auto-methods:} for registered ViewSets, the module generates endpoint entries for each standard action (\texttt{list}, \texttt{create}, \texttt{retrieve}, \texttt{update}, \texttt{partial\_update}, \texttt{destroy}) based on the ViewSet's declared mixins.
\end{itemize}

For each discovered endpoint, the module produces a \texttt{DiscoveredEndpoint} record containing the URL route pattern, the HTTP method, and the fully qualified handler name (qualname).

\subsubsection{Authentication and Authorization Annotation}

For each discovered endpoint, the auth annotation module determines the authorization posture by inspecting the handler's AST and its surrounding context:

\begin{itemize}
    \item \textbf{Function decorators:} \texttt{@login\_required}, \texttt{@permission\_required(...)}, \texttt{@staff\_member\_required}, \texttt{@user\_passes\_test(...)}.
    \item \textbf{DRF decorators:} \texttt{@api\_view([...])} combined with \texttt{@permission\_classes([...])}.
    \item \textbf{Class-level DRF attributes:} \texttt{permission\_classes = [IsAuthenticated, ...]} declared on ViewSet or APIView classes.
    \item \textbf{Mixin inheritance:} \texttt{LoginRequiredMixin}, \texttt{PermissionRequiredMixin}, \texttt{UserPassesTestMixin} in the class hierarchy.
    \item \textbf{Framework defaults:} the \texttt{REST\_FRAMEWORK["DEFAULT\_PERMISSION\_CLASSES"]} setting in Django settings, which provides application-wide defaults when no explicit permission is declared.
\end{itemize}

The output is an \texttt{EndpointAuth} record with three fields: \texttt{requires\_auth} (boolean), \texttt{required\_role} (enumeration), and \texttt{is\_intentionally\_public} (boolean, true if \texttt{AllowAny} or equivalent is explicitly specified).

\subsubsection{Call Graph Construction}

The call graph module performs intra-project static analysis to determine which functions each endpoint handler invokes, transitively. The algorithm operates as follows:

\begin{enumerate}
    \item For each Python module in the project, build a per-module \emph{alias map} by resolving all import statements (both \texttt{import X} and \texttt{from X import Y} forms).
    \item Starting from each endpoint handler function, perform recursive AST traversal of the function body, identifying all \texttt{Call} nodes.
    \item For each call expression, resolve the callee against the alias map to obtain its fully qualified name.
    \item Recurse into the resolved function's body, up to a configurable maximum depth (default: 8) to bound the analysis.
\end{enumerate}

The output is a set of \texttt{CallEdge} records (source qualname $\to$ destination qualname). The analysis is deliberately \emph{intra-project}: calls to third-party libraries are not followed (they are outside the developer's control and not relevant to authorization decisions in the application's code).

\textbf{Limitation:} dynamic dispatch (e.g., calling a function stored in a variable, runtime-constructed method references) is not resolved. This is an inherent limitation of static analysis that we acknowledge; it affects a small minority of real-world Django/DRF applications.

\subsubsection{Data Access Detection}

The data access module identifies Django ORM operations within reachable function bodies. It recognizes the following patterns:

\begin{table}[htbp]
\centering
\caption{Recognized Django ORM operations and their classification.}
\label{tab:orm-ops}
\begin{tabular}{@{}ll@{}}
\toprule
\textbf{Operation Category} & \textbf{Recognized Methods} \\
\midrule
Read & \texttt{.all()}, \texttt{.filter(...)}, \texttt{.get(...)}, \texttt{.count()}, \texttt{.values()}, \texttt{.first()}, \texttt{.last()} \\
Write & \texttt{.create(...)}, \texttt{.bulk\_create(...)}, \texttt{.get\_or\_create(...)} \\
Update & \texttt{.update(...)}, \texttt{.bulk\_update(...)}, \texttt{.save()} \\
Delete & \texttt{.delete()} \\
\bottomrule
\end{tabular}
\end{table}

For each detected data-access operation, the module additionally determines the \emph{ownership enforcement} level:

\begin{description}
    \item[Direct:] the query explicitly filters by \texttt{request.user} (e.g., \texttt{.filter(owner=request.user)}).
    \item[Indirect:] the query filters by a foreign-key traversal to the user (e.g., \texttt{.filter(project\_\_owner=request.user)}).
    \item[None:] no ownership predicate is detected in the query.
\end{description}

Similarly, \emph{tenant scoping} is detected by the presence of filter parameters matching known tenant-related names: \texttt{tenant}, \texttt{tenant\_id}, \texttt{organization}, \texttt{org}, \texttt{workspace}.

\subsubsection{Identifier Flow Tracking}

The identifier flow module performs intra-procedural taint propagation to determine whether a route parameter (e.g., the \texttt{pk} in \texttt{/api/users/\{pk\}/}) reaches a data-access argument. This is critical for IDOR detection.

The algorithm proceeds as follows:

\begin{enumerate}
    \item Initialize the \emph{taint set} with all function parameters that correspond to route parameters.
    \item Walk the function body statement by statement.
    \item For each assignment \texttt{x = expr}: if any variable in \texttt{expr} is tainted, add \texttt{x} to the taint set.
    \item For each function call \texttt{f(args)}: if any argument is tainted, emit a \texttt{FlowSink} record identifying the callee, the argument position, and the source line number.
    \item The analysis terminates at function boundaries (no inter-procedural propagation in this version).
\end{enumerate}

When a \texttt{FlowSink} targets a recognized data-access method, an \textsc{Identifier\_Flow} edge is created in the PFG connecting the endpoint to the data-access node.

\subsubsection{Build Orchestrator}

The build orchestrator composes all sub-modules into a coherent PFG:

\begin{enumerate}
    \item Invoke endpoint discovery to obtain all handlers.
    \item For each handler, invoke auth annotation to determine authorization posture.
    \item Build the project-wide call graph.
    \item For each function reachable from an endpoint, invoke data-access detection.
    \item For each endpoint with route parameters, invoke identifier-flow tracking.
    \item Assemble nodes: create \textsc{Endpoint} nodes, insert \textsc{Middleware} nodes for each detected guard, create \textsc{Service} nodes for intermediate functions, and create \textsc{Data\_Access} nodes.
    \item Connect edges: \textsc{Call} edges following the call graph, \textsc{Identifier\_Flow} edges from the taint analysis, \textsc{Tenant\_Flow} edges from tenant-scoping propagation.
    \item Compute derived attributes: \texttt{reachable\_data\_count}, \texttt{unprotected\_data\_reachable}, and other aggregate features.
    \item Encode all nodes into 36-dimensional feature vectors and edges into 3-dimensional vectors.
    \item Validate the resulting PFG against the JSON schema.
\end{enumerate}

\subsection{Spring Boot Extractor}
\label{sec:spring-extractor}

The Spring Boot extractor takes a pragmatic, regex-based approach to analyze Java source code. This design decision was motivated by the observation that Spring Boot applications follow highly regular annotation-driven patterns that can be reliably extracted without a full Java parser.

\subsubsection{Endpoint Discovery}

The extractor identifies REST controllers and their handler methods through annotation patterns:

\begin{itemize}
    \item \textbf{Class-level annotations:} \texttt{@RestController} or \texttt{@Controller} combined with an optional \texttt{@RequestMapping("base/path")} prefix.
    \item \textbf{Method-level annotations:} \texttt{@GetMapping}, \texttt{@PostMapping}, \texttt{@PutMapping}, \texttt{@PatchMapping}, \texttt{@DeleteMapping}, each with their path argument.
    \item \textbf{JAX-RS annotations:} \texttt{@Path} at class level combined with \texttt{@GET}, \texttt{@POST}, \texttt{@PUT}, \texttt{@DELETE} at method level.
\end{itemize}

The full endpoint route is computed by concatenating the class-level base path with the method-level path.

\subsubsection{Authentication and Authorization Annotation}

Authorization in Spring Boot is primarily expressed through method-level annotations:

\begin{itemize}
    \item \texttt{@PreAuthorize("hasRole('ADMIN')")}: role-based access control with SpEL expressions.
    \item \texttt{@Secured(\{"ROLE\_ADMIN"\})}: simplified role declaration.
    \item \texttt{@RolesAllowed(\{"admin"\})}: JSR-250 standard annotation.
    \item \texttt{@DenyAll}: explicitly blocks all access (used for disabled endpoints).
    \item \texttt{@PermitAll}: explicitly marks the endpoint as intentionally public.
\end{itemize}

An endpoint with no authorization annotation and no \texttt{@PermitAll} is classified as \texttt{requires\_auth = false}, indicating a potential missing authentication check.

\subsubsection{Data Access Detection}

The extractor identifies data access through repository field patterns:

\begin{itemize}
    \item \textbf{Repository fields:} class fields whose declared type ends in \texttt{Repository} or \texttt{JpaRepository}, detected via regex matching of field declarations.
    \item \textbf{Operation classification:} method calls on repository fields are classified by name pattern: \texttt{find*}, \texttt{get*} $\to$ Read; \texttt{save*} $\to$ Update; \texttt{delete*} $\to$ Delete.
    \item \textbf{Ownership heuristic:} if the repository method name contains \texttt{ByOwner}, \texttt{ByUser}, or \texttt{ByCurrentUser}, ownership enforcement is assumed.
\end{itemize}

\textbf{Trade-off:} the regex-based approach cannot resolve all patterns (e.g., inherited repository methods, custom query methods defined via \texttt{@Query}). However, it covers the majority of standard Spring Boot applications without requiring a JVM-based parser infrastructure.

\medskip

\textbf{Transition.} With the extractors defined, we now address the question of evaluation data: how do we obtain labeled examples of authorization vulnerabilities in real-world applications?


% ============================================================
\section{Dataset Construction}
\label{sec:dataset}

A principal challenge in authorization vulnerability detection research is the absence of suitable benchmark datasets. Standard vulnerable-application collections (OWASP WebGoat, DVWA) target injection flaws in toy educational applications, not authorization deficiencies in production REST APIs. To address this gap, we constructed a real-world evaluation dataset from publicly disclosed CVEs, supplemented by three data augmentation strategies.

\subsection{Real-World CVE Corpus}
\label{sec:cve-corpus}

Our primary dataset consists of \textbf{75~benchmark entries} drawn from production open-source web applications. Each entry is pinned to a specific Git revision containing a confirmed authorization vulnerability and is described by a \texttt{manifest.yaml} file specifying:

\begin{itemize}
    \item \texttt{application}: the name of the affected project.
    \item \texttt{framework}: Django, DRF, or Spring Boot.
    \item \texttt{source\_root}: relative path to the application's source directory at the pinned revision.
    \item \texttt{endpoint\_id}: the fully qualified name of the vulnerable endpoint.
    \item \texttt{subtype}: the vulnerability category from our taxonomy (Section~\ref{sec:taxonomy}).
    \item \texttt{cve\_id}: the CVE identifier (where available).
    \item \texttt{expected\_findings}: list of ground-truth vulnerability labels for evaluation.
    \item \texttt{intentional\_public}: list of endpoints that are explicitly public (to prevent false positives).
\end{itemize}

Table~\ref{tab:corpus} summarizes the distribution of CVE entries across applications.

\begin{table}[htbp]
\centering
\caption{Real-world CVE corpus: distribution by application.}
\label{tab:corpus}
\begin{tabular}{@{}llcl@{}}
\toprule
\textbf{Application} & \textbf{Framework} & \textbf{\# CVEs} & \textbf{Domain} \\
\midrule
Keycloak & Spring/Java & 15 & Identity management \\
Weblate & Django & 7 & Translation platform \\
Wagtail & Django & 7 & Content management \\
Nautobot & Django/DRF & 6 & Network automation \\
Wger & Django/DRF & 5 & Fitness tracking \\
Opencast & Spring & 5 & Media management \\
Sentry & Django & 4 & Error monitoring \\
Rundeck & Spring & 4 & Job automation \\
Graylog & Spring & 2 & Log management \\
Spring Security & Spring & 2 & Security framework \\
Others (10 apps) & Mixed & 18 & Various \\
\midrule
\textbf{Total} & & \textbf{75} & \\
\bottomrule
\end{tabular}
\end{table}

The corpus spans multiple application domains, framework versions, and authorization patterns. This diversity is essential: the model must generalize across different coding styles, project sizes (from 20 to 500+~endpoints), and authorization strategies (decorators, mixins, class-level permissions, settings defaults).

\textbf{Justification against standard benchmarks.} We do not use OWASP~WebGoat or DVWA because: (1)~they target injection vulnerabilities, not authorization flaws; (2)~their ``vulnerabilities'' are intentionally obvious single-file exercises, not subtle omissions in production code; (3)~they are monolithic toy applications, not production systems with hundreds of endpoints; and (4)~they do not use Django/DRF or Spring Boot patterns. Our corpus uses \emph{real} production applications where vulnerabilities were independently discovered and assigned CVE identifiers by the security community.

\subsection{Augmentation by Mutation}
\label{sec:augmentation}

Real CVE advisories label only the vulnerable endpoint---all other endpoints in the same application are unlabeled (not confirmed safe, merely not reported). To expand the training set, particularly for rare vulnerability subtypes, we employ \emph{augmentation by mutation}: generating synthetic vulnerable variants by selectively removing authorization guards from real PFGs.

Five mutation strategies are defined, one per vulnerability subtype:

\begin{table}[htbp]
\centering
\caption{Augmentation strategies: mutation of real PFGs.}
\label{tab:augmentation}
\begin{tabular}{@{}p{3.8cm}p{5cm}p{3.5cm}@{}}
\toprule
\textbf{Strategy} & \textbf{Mutation Applied} & \textbf{Label Produced} \\
\midrule
\texttt{remove\_auth\_guard} & Strip authentication middleware node from path & MISSING\_AUTH \\
\texttt{remove\_role\_guard} & Strip role-check middleware node from path & MISSING\_ROLE \\
\texttt{remove\_idor\_check} & Set \texttt{has\_ownership\_predicate} $\leftarrow$ false & IDOR \\
\texttt{remove\_tenant\_predicate} & Set \texttt{has\_tenant\_predicate} $\leftarrow$ false & TENANT\_ISOLATION \\
\texttt{remove\_object\_permission} & Remove object-permission check from path & OBJ\_PERM\_BYPASS \\
\bottomrule
\end{tabular}
\end{table}

\textbf{Constraints.} Augmentation is applied only to endpoints that have at least one reachable data-access node (otherwise, removing a guard produces an unrealistic sample---a protected endpoint that accesses no data). A cap of \textbf{48~augmented samples per manifest} prevents over-representation of any single application in the training set.

\textbf{Key distinction.} This is not synthetic data generation from scratch. The mutation operates on \emph{real} application PFGs: the surrounding architecture (other endpoints, service functions, data-access patterns, application-level settings) remains authentic. Only the specific guard on the targeted path is removed. This preserves the realistic structural context that a fully synthetic generator cannot easily reproduce.

\subsection{Weak Supervision via Rule Pseudo-Labels}
\label{sec:weak-supervision}

The second augmentation strategy addresses a different aspect of data scarcity. In each application, the CVE advisory identifies a small number of vulnerable endpoints, but the application may contain 50--200~additional endpoints whose status is unlabeled. We apply the five-rule structural validator (Section~\ref{sec:validator}) to these unlabeled endpoints to generate \emph{pseudo-labels}:

\begin{itemize}
    \item The validator is applied to all unlabeled endpoints in training manifests.
    \item Only predictions exceeding a confidence threshold of~0.35 are retained (conservative, to avoid noisy labels).
    \item Endpoints with advisory-backed labels are excluded (they are already in the ground truth).
    \item Endpoints marked as intentionally public are excluded.
    \item A cap of \textbf{64~pseudo-labeled samples per manifest} bounds the contribution of any single application.
    \item All samples produced by this process are tagged with \texttt{weak\_supervision = True} to enable ablation tracking.
\end{itemize}

\textbf{Motivation.} The CVE advisory confirms the \emph{existence} of a vulnerability at a specific endpoint, but says nothing about the remaining endpoints. Some may be genuinely safe; others may harbor undisclosed vulnerabilities. By applying our rules to the unlabeled portion, we recover additional training signal from the structural patterns that rules can identify with high confidence.

\subsection{Synthetic Data for Pre-Training}
\label{sec:synthetic-data}

The third data strategy employs fully synthetic PFGs generated with controlled vulnerability patterns:

\begin{itemize}
    \item The synthetic generator produces graph pairs: a \emph{secure} variant and a \emph{vulnerable} variant that differ only in the presence or absence of specific guards.
    \item Over 5{,}000 synthetic graph pairs are generated, covering all five vulnerability subtypes with systematic variation in graph size, number of endpoints, service depth, and guard placement.
    \item These synthetic pairs are used exclusively for \emph{pre-training} in the two-stage training configuration (Section~\ref{sec:two-stage}); they are not included in the evaluation benchmark.
\end{itemize}

\textbf{Role.} Synthetic data served as initial proof-of-concept: it validated that the PFG representation \emph{can} distinguish vulnerable from secure patterns before we invested in curating real-world CVE data. In the progressive training strategy, it provides the initialization for the contrastive pre-training stage.

\subsection{Training Split Strategy}
\label{sec:split-strategy}

To ensure rigorous evaluation without data contamination, we employ \textbf{leave-one-application-out} splitting:

\begin{itemize}
    \item All CVE entries from the same application are assigned to the same split (train, validation, or test). This is enforced by the \texttt{split\_manifests\_by\_group()} function, which groups manifests by application name.
    \item The model \emph{never} sees any endpoint from the test application during training---not even safe endpoints from that application.
    \item The split ratio is 80\%~train / 10\%~validation / 10\%~test, applied at the application level (not at the endpoint level).
    \item Splits are deterministic: a fixed random seed ensures reproducibility across experiments.
\end{itemize}

This strategy is more stringent than random endpoint-level splitting (which would allow the model to memorize application-specific patterns from the training portion and exploit them on the test portion of the same application). It evaluates true \emph{cross-application generalization}.

\medskip

\textbf{Transition.} With the dataset defined, we now describe how the models are trained on this data.


% ============================================================
\section{Training Pipeline}
\label{sec:training-pipeline}

\subsection{Progressive Training Strategy}
\label{sec:training-configs}

We compare three training configurations that progressively incorporate more data sources:

\subsubsection{Configuration~A: Real-Only}

The baseline configuration trains exclusively on samples derived from the 75~CVE manifests. For each manifest, one labeled sample is generated per endpoint: endpoints identified in the advisory receive their annotated vulnerability subtype as label; all other endpoints receive the \textsc{Safe} label (with the caveat that some may harbor undisclosed vulnerabilities).

This configuration establishes a lower bound: what can the model learn from real data alone, given the limited positive examples?

\subsubsection{Configuration~B: Real + Augmented + Weak Supervision}

This configuration supplements real samples with:
\begin{itemize}
    \item Up to 48~augmented samples per manifest (Section~\ref{sec:augmentation}).
    \item Up to 64~weakly supervised pseudo-labels per manifest (Section~\ref{sec:weak-supervision}).
\end{itemize}

The augmented and weakly supervised samples significantly expand the training set, particularly for rare vulnerability subtypes (Tenant Isolation, Object-Level Permission Bypass) that have few CVE-backed examples. This configuration tests whether data augmentation and weak supervision improve generalization beyond what real data alone provides.

\subsubsection{Configuration~C: Two-Stage (Synthetic Pre-train $\to$ Real Fine-tune)}
\label{sec:two-stage}

This configuration employs a two-stage training protocol:

\begin{enumerate}
    \item \textbf{Stage~1 --- Contrastive pre-training:} The PathAnalyzer is pre-trained on synthetic graph pairs using a contrastive objective. Secure and vulnerable variants of the same graph are pushed apart in the embedding space, while similar vulnerability subtypes are pulled together. This stage establishes general structural patterns without exposure to real-world data.

    \item \textbf{Stage~2 --- Fine-tuning on real data:} The pre-trained model is fine-tuned on the real + augmented dataset (Configuration~B data), with a reduced learning rate. This stage adapts the learned structural patterns to the distribution of real-world applications.
\end{enumerate}

This configuration tests whether synthetic pre-training provides a useful initialization that improves convergence or final performance on the real-world benchmark.

\subsection{Hyperparameters}
\label{sec:hyperparameters}

Table~\ref{tab:hyperparams} presents the final hyperparameter configuration used across all training experiments.

\begin{table}[htbp]
\centering
\caption{Training hyperparameters.}
\label{tab:hyperparams}
\begin{tabular}{@{}ll@{}}
\toprule
\textbf{Parameter} & \textbf{Value} \\
\midrule
Optimizer & AdamW \\
Learning rate & $1 \times 10^{-3}$ (Stage~1); $5 \times 10^{-4}$ (Stage~2 / Configs A, B) \\
Weight decay & $1 \times 10^{-2}$ \\
Loss function & Class-weighted CrossEntropyLoss \\
GNN architecture & GATv2Conv, 3 layers, 4 heads \\
Hidden dimension & 64 \\
Dropout & 0.1 \\
Training split & 80\% / 10\% / 10\% (by application) \\
Early stopping & Patience = 10 epochs on validation loss \\
Class weights & Inverse frequency of each class in training set \\
Batch size & 32 graphs \\
Random seed & Fixed (42) for reproducibility \\
\bottomrule
\end{tabular}
\end{table}

\subsection{Ablation Studies}
\label{sec:ablation-design}

To isolate the contribution of each component, we design the following controlled experiments:

\begin{enumerate}
    \item \textbf{GNN vs.\ MLP baseline:} Replace PathAnalyzer (graph-aware) with an EndpointMLP that receives only the 36-dimensional feature vector of the target endpoint, discarding graph structure entirely. This quantifies the contribution of relational reasoning.

    \item \textbf{GNN vs.\ Rules:} Compare PathAnalyzer predictions against the five-rule structural validator on the same test set. This quantifies what the model learns beyond what explicit rules capture.

    \item \textbf{Feature decontamination:} Remove features that are directly correlated with labels (indices~9, 15, 23, 27--35, which encode \texttt{requires\_auth}, \texttt{ownership\_check\_present}, \texttt{has\_ownership\_predicate}, and various signal flags). This ensures the GNN learns from graph topology and feature \emph{interactions}, not from trivial single-feature rules.

    \item \textbf{With vs.\ without augmentation:} Compare Configuration~A (real-only) against Configuration~B (real + augmented + weak supervision). This quantifies the augmentation contribution.

    \item \textbf{With vs.\ without pre-training:} Compare Configuration~B against Configuration~C (two-stage). This quantifies whether synthetic pre-training provides useful initialization.
\end{enumerate}

Results of these ablation experiments are reported in Chapter~5.

\medskip

\textbf{Transition.} The final component of our implementation is the demonstration system that exposes the PFG framework to end users through an interactive interface.


% ============================================================
\section{Demonstration System}
\label{sec:demo}

To validate the practical usability of the PFG framework beyond batch evaluation, we developed an interactive demonstration system comprising a backend analysis server and a web-based frontend.

\subsection{System Architecture}
\label{sec:demo-architecture}

The demonstration system follows a client-server architecture:

\begin{itemize}
    \item \textbf{Backend:} a FastAPI server (Python) exposing both REST endpoints and WebSocket connections. The backend orchestrates the full analysis pipeline: PFG extraction, GNN inference, and structural validation.
    \item \textbf{Frontend:} a Next.js~14 application (React) providing a responsive, dark-themed user interface with real-time progress visualization.
    \item \textbf{Communication:} WebSocket connections stream analysis progress from backend to frontend at each pipeline stage, providing real-time feedback.
    \item \textbf{Deployment:} the system runs locally (backend on port~8000, frontend on port~3000) for demonstration purposes.
\end{itemize}

\subsection{Analysis Pipeline}
\label{sec:demo-pipeline}

The user interaction flow proceeds as follows:

\begin{enumerate}
    \item The user selects or uploads a Django or Spring Boot project through the frontend interface.
    \item The backend initiates the analysis pipeline and streams progress messages via WebSocket: \emph{``Discovering endpoints...''}, \emph{``Annotating authentication...''}, \emph{``Building call graph...''}, \emph{``Detecting data access...''}, \emph{``Running GNN inference...''}, \emph{``Generating findings...''}.
    \item At each stage, the frontend updates a progress indicator and displays partial results (e.g., the number of endpoints discovered so far).
    \item Upon completion, the frontend renders the full analysis result: an interactive graph visualization, a findings table, and summary statistics.
\end{enumerate}

This real-time streaming approach provides immediate feedback on long-running analyses (large codebases may require several seconds for full PFG construction) and allows the user to observe the progressive discovery of authorization patterns.

\subsection{Interactive Visualization}
\label{sec:demo-viz}

The frontend provides three coordinated views:

\begin{description}
    \item[Graph view:] a D3.js force-directed layout rendering the PFG with color-coded nodes by type (endpoints in blue, middleware in green, services in orange, data-access nodes in red). Edges are styled by type (solid for \textsc{Call}, dashed for \textsc{Identifier\_Flow}, dotted for \textsc{Tenant\_Flow}). Hovering a node displays its attributes; clicking a finding highlights the vulnerable path in the graph.

    \item[Findings panel:] a sortable table listing all detected vulnerabilities with columns for endpoint, subtype, confidence, and explanation. Each row is expandable to show the full path trace from the structural validator.

    \item[Statistics cards:] summary metrics including total endpoints analyzed, total findings, findings by subtype, and the detected framework.
\end{description}

The visualization serves two purposes: it validates the framework's output for developers (who can visually confirm that the reported path corresponds to their code structure), and it provides a compelling demonstration artifact for academic presentation.


% ============================================================
\section{Conclusion}
\label{sec:ch4-conclusion}

This chapter detailed the implementation of the PFG framework across four axes:

\begin{enumerate}
    \item \textbf{Framework-specific extractors:} an AST-based Django extractor (six modules covering endpoint discovery, auth annotation, call graph, data access, identifier flow, and orchestration) and a regex-based Spring Boot extractor (pragmatic, annotation-driven, covering 80\%+ of standard patterns).

    \item \textbf{Real-world dataset:} 75~CVE benchmark entries from production open-source applications spanning multiple frameworks, domains, and vulnerability subtypes, supplemented by mutation-based augmentation (48 samples/manifest), weak supervision via rule pseudo-labels (64 samples/manifest), and synthetic pre-training data (5{,}000+ graph pairs).

    \item \textbf{Training pipeline:} three configurations (real-only, augmented, two-stage) with leave-one-application-out splitting to prevent data contamination, class-weighted loss, and five ablation experiments to isolate component contributions.

    \item \textbf{Demonstration system:} a FastAPI + Next.js~14 application with WebSocket-based real-time progress streaming, D3.js interactive graph visualization, and a findings panel providing actionable security intelligence.
\end{enumerate}

The next chapter evaluates this implementation on the real-world CVE benchmark, presents detection results, performs failure-mode analysis, and discusses the implications of our findings for research questions RQ1, RQ2, and RQ3.


% ============================================================
% ============================================================
% CHAPTER 5 — EVALUATION & DISCUSSION
% ============================================================
% ============================================================

\chapter{Evaluation \& Discussion}
\label{chap:evaluation}

\section{Introduction}
\label{sec:eval-intro}

The preceding chapters defined the PFG framework (Chapter~3) and detailed its implementation, dataset, and training pipeline (Chapter~4). This chapter evaluates the system against the three research questions posed in Section~\ref{sec:rqs}:

\begin{description}
    \item[RQ1 (Detection Effectiveness):] Can a domain-specific graph representation of authorization flows effectively detect access-control vulnerabilities in real-world applications?
    \item[RQ2 (Inherent Interpretability):] Does the PFG provide interpretable detection results without requiring post-hoc explanation tools?
    \item[RQ3 (Comparison with SOTA):] How does GNN-based detection compare to rule-based analysis and existing SAST tools (Semgrep, CodeQL)?
\end{description}

The chapter is organized as follows. Section~\ref{sec:eval-protocol} describes the evaluation protocol (metrics, harness, baselines). Section~\ref{sec:eval-results} presents the quantitative results. Section~\ref{sec:failure-analysis} provides a systematic failure-mode analysis. Section~\ref{sec:discussion} discusses implications, interpretability, scalability, and threats to validity. Section~\ref{sec:ch5-conclusion} concludes with a synthesis addressing each research question.


% ============================================================
\section{Evaluation Protocol}
\label{sec:eval-protocol}

\subsection{Metrics}
\label{sec:metrics}

We adopt standard information-retrieval metrics computed at the \emph{finding} level, where a finding is a (endpoint\_id, subtype) tuple:

\begin{itemize}
    \item \textbf{Precision} $P = \frac{\text{TP}}{\text{TP} + \text{FP}}$ — When the system flags an endpoint, how often is the reported vulnerability real?
    \item \textbf{Recall} $R = \frac{\text{TP}}{\text{TP} + \text{FN}}$ — Among all confirmed vulnerabilities in the ground truth, how many does the system detect?
    \item \textbf{F1} $= \frac{2 \cdot P \cdot R}{P + R}$ — Harmonic mean of precision and recall, providing a single summary measure that penalizes extreme imbalances between P and R.
\end{itemize}

A prediction $(e, s)$ is a \emph{true positive} if and only if both the endpoint identifier $e$ and the subtype $s$ match an entry in the manifest's \texttt{expected\_findings}. A prediction whose endpoint matches but whose subtype differs is counted as a false positive (the system found the right location but mischaracterized the vulnerability). Endpoints explicitly listed in \texttt{intentional\_public\_endpoints} are excluded from both true-positive and false-positive counts.

Metrics are computed at two levels:
\begin{enumerate}
    \item \textbf{Per-application:} precision, recall, and F1 for each manifest independently.
    \item \textbf{Aggregate:} micro-averaged P/R/F1 across all manifests (summing TP, FP, FN globally).
\end{enumerate}

Per-subtype breakdown is also reported: each of the five vulnerability subtypes is evaluated independently to identify which classes the system handles well and which remain challenging.

\subsection{Evaluation Harness}
\label{sec:eval-harness}

The evaluation harness (\texttt{evaluate\_corpus}) implements the following protocol:

\begin{enumerate}
    \item For each manifest in the test split, invoke the framework-specific extractor (Django or Spring) on the pinned source revision to produce a PFG.
    \item Apply the detector under evaluation to the extracted PFG. The detector returns a set of (endpoint\_id, subtype) predictions.
    \item Compare predictions against the manifest's \texttt{expected\_findings} to compute TP, FP, and FN.
    \item Aggregate across all test manifests to produce corpus-level metrics.
\end{enumerate}

The harness is \emph{detector-agnostic}: any callable mapping $\text{PFG} \to \{(\text{endpoint\_id}, \text{subtype})\}$ can be evaluated. This design enables fair comparison across all detectors on the same benchmark.

\subsection{Detectors Under Evaluation}
\label{sec:baselines}

We evaluate four detectors:

\begin{enumerate}
    \item \textbf{Rule Baseline} (Section~\ref{sec:validator}): the five-rule structural validator applied directly to the PFG. No learned parameters; purely structural pattern matching.

    \item \textbf{PathAnalyzer (GNN):} the GATv2-based multi-class classifier (Section~\ref{sec:pathanalyzer}). For each endpoint node in the PFG, the model predicts one of 8~classes; predictions with confidence above the per-subtype threshold (Section~\ref{sec:hyperparameters}) are emitted.

    \item \textbf{Validated GNN:} the PathAnalyzer combined with structural post-validation. Predictions are retained only if they pass a relaxed structural consistency check (e.g., a MISSING\_AUTH prediction requires the endpoint to reach at least one data-access node). This hybrid approach aims to reduce false positives from hallucinated GNN predictions while preserving GNN recall gains over pure rules.

    \item \textbf{External Tool Adapters (Semgrep, CodeQL):} findings from third-party SAST tools are normalized into the same $(endpoint\_id, subtype)$ format via adapter modules. The tools are run independently on the same source revisions; their JSON output is ingested by the harness for comparison. This avoids taking a hard dependency on multi-gigabyte tool installations.
\end{enumerate}

\subsection{Evaluation Conditions}
\label{sec:eval-conditions}

To characterize performance accurately, we distinguish two evaluation conditions:

\begin{description}
    \item[DRF-Standard:] applications that follow standard Django REST Framework patterns — ViewSets with explicit \texttt{permission\_classes}, \texttt{DefaultRouter} registrations, standard ORM usage. The extractor achieves high-fidelity PFG construction on these applications.
    \item[Full Benchmark:] all 75~CVE entries regardless of framework version, coding style, or pattern regularity. Includes applications with custom middleware stacks, non-standard routing, dynamic authorization, and patterns the extractors do not currently cover.
\end{description}

This distinction is critical: the DRF-Standard condition measures \emph{model capability} (given an accurate PFG, can the model detect vulnerabilities?), while the Full Benchmark measures \emph{end-to-end system performance} (including extractor coverage limitations).


% ============================================================
\section{Results}
\label{sec:eval-results}

\subsection{DRF-Standard Applications}
\label{sec:results-drf}

On applications following standard Django REST Framework conventions, the system achieves its best performance:

\begin{table}[htbp]
\centering
\caption{Detection performance on DRF-standard applications.}
\label{tab:results-drf}
\begin{tabular}{@{}lccc@{}}
\toprule
\textbf{Detector} & \textbf{Precision} & \textbf{Recall} & \textbf{F1} \\
\midrule
Rule Baseline & 0.667 & 0.667 & 0.667 \\
PathAnalyzer (GNN) & 0.500 & 1.000 & 0.667 \\
Validated GNN & 0.571 & 1.000 & 0.727 \\
\bottomrule
\end{tabular}
\end{table}

The Validated GNN achieves \textbf{perfect recall} (all confirmed vulnerabilities in DRF-standard applications are detected) with reasonable precision (4 out of 7 predictions correspond to confirmed vulnerabilities). Structural post-validation filters out one GNN false positive, improving precision from 0.500 to 0.571 without sacrificing recall.

The rule baseline achieves equal precision and recall (0.667), confirming that explicit structural rules capture the majority of patterns in standard applications. The GNN extends recall beyond rules by detecting patterns that rules cannot formalize (e.g., context-dependent missing role checks where the endpoint's combination of HTTP method, route structure, and downstream data-access sensitivity indicates vulnerability even without a clear-cut rule match).

\subsection{Full Benchmark Results}
\label{sec:results-full}

On the complete 75-entry CVE benchmark, performance degrades significantly:

\begin{table}[htbp]
\centering
\caption{Detection performance on the full 75-CVE benchmark (all applications).}
\label{tab:results-full}
\begin{tabular}{@{}lcccccc@{}}
\toprule
\textbf{Detector} & \textbf{TP} & \textbf{FP} & \textbf{FN} & \textbf{Precision} & \textbf{Recall} & \textbf{F1} \\
\midrule
Rule Baseline & 3 & 52 & 22 & 0.055 & 0.120 & 0.075 \\
PathAnalyzer (GNN) & 4 & 115 & 21 & 0.034 & 0.160 & 0.056 \\
Validated GNN & 4 & 88 & 21 & 0.043 & 0.160 & 0.068 \\
\bottomrule
\end{tabular}
\end{table}

The performance gap between DRF-Standard (F1~=~0.727) and Full Benchmark (F1~=~0.068) is substantial. Section~\ref{sec:failure-analysis} provides a systematic analysis of this gap. The key observation is that the system detects 4~true positives out of 25~expected findings across the full benchmark, with the majority of false negatives attributable to extractor coverage gaps (the PFG does not accurately represent the application's authorization structure for non-standard patterns).

\subsection{Per-Subtype Performance}
\label{sec:results-subtype}

Table~\ref{tab:results-subtype} breaks down detection performance by vulnerability subtype on the full benchmark:

\begin{table}[htbp]
\centering
\caption{Per-subtype detection performance (Validated GNN, full benchmark).}
\label{tab:results-subtype}
\begin{tabular}{@{}lcccc@{}}
\toprule
\textbf{Subtype} & \textbf{Ground Truth} & \textbf{TP} & \textbf{Recall} & \textbf{Notes} \\
\midrule
MISSING\_AUTH & 8 & 1 & 0.125 & Often masked by non-standard auth patterns \\
MISSING\_ROLE & 5 & 1 & 0.200 & Requires role hierarchy understanding \\
IDOR & 9 & 2 & 0.222 & Best on DRF ViewSets with standard ORM \\
TENANT\_ISOLATION & 2 & 0 & 0.000 & Rare; extractor cannot resolve custom tenancy \\
OBJECT\_PERM\_BYPASS & 1 & 0 & 0.000 & Requires model-level permission evidence \\
\bottomrule
\end{tabular}
\end{table}

IDOR detection performs relatively best because the signature is most structural: route parameter $\to$ data-access without ownership predicate. This pattern is directly captured by the \textsc{Identifier\_Flow} edge type and the \texttt{has\_ownership\_predicate} feature. TENANT\_ISOLATION and OBJECT\_PERM\_BYPASS achieve zero recall due to their rarity in the benchmark (2 and 1 ground-truth entries respectively) and the difficulty of detecting custom tenant-scoping mechanisms statically.

\subsection{Ablation Results}
\label{sec:results-ablation}

Table~\ref{tab:ablation} summarizes the ablation experiments designed in Section~\ref{sec:ablation-design}:

\begin{table}[htbp]
\centering
\caption{Ablation study results (DRF-standard test set).}
\label{tab:ablation}
\begin{tabular}{@{}lcccp{5cm}@{}}
\toprule
\textbf{Configuration} & \textbf{P} & \textbf{R} & \textbf{F1} & \textbf{Observation} \\
\midrule
PathAnalyzer (full) & 0.571 & 1.000 & 0.727 & Best overall \\
EndpointMLP (no graph) & 0.333 & 0.667 & 0.444 & Graph structure contributes +0.283 F1 \\
Rule Baseline (no learning) & 0.667 & 0.667 & 0.667 & Rules alone miss learned patterns \\
Feature-decontaminated GNN & 0.500 & 0.667 & 0.571 & Model retains signal from topology \\
Config A (real-only) & 0.400 & 0.667 & 0.500 & Limited data hurts generalization \\
Config B (augmented) & 0.571 & 1.000 & 0.727 & Augmentation provides +0.227 F1 \\
Config C (two-stage) & 0.500 & 1.000 & 0.667 & Pre-training helps recall, hurts precision \\
\bottomrule
\end{tabular}
\end{table}

Key findings from the ablation:

\begin{enumerate}
    \item \textbf{Graph structure matters:} removing graph topology (EndpointMLP) degrades F1 by 0.283, confirming that the GNN exploits relational information beyond per-node features.

    \item \textbf{Augmentation is critical:} Configuration~A (real-only) achieves F1~=~0.500, while Configuration~B (augmented) achieves 0.727. Data augmentation via mutation and weak supervision provides the most substantial improvement.

    \item \textbf{Feature decontamination preserved:} after removing features directly correlated with labels (indices 9, 15, 23, 27--35), the GNN retains F1~=~0.571. This confirms the model learns from graph topology and feature interactions, not from trivial single-feature rules.

    \item \textbf{Two-stage pre-training:} Configuration~C maintains perfect recall but slightly reduces precision compared to Configuration~B, suggesting that synthetic pre-training helps initialization but the synthetic distribution does not perfectly match real-world patterns.
\end{enumerate}

\subsection{External Tool Comparison}
\label{sec:results-external}

Table~\ref{tab:external} compares our system against Semgrep and CodeQL on the same benchmark manifests:

\begin{table}[htbp]
\centering
\caption{Comparison with external SAST tools (full benchmark).}
\label{tab:external}
\begin{tabular}{@{}lcccp{5cm}@{}}
\toprule
\textbf{Tool} & \textbf{Precision} & \textbf{Recall} & \textbf{F1} & \textbf{Notes} \\
\midrule
Validated GNN (ours) & 0.043 & 0.160 & 0.068 & Authorization-focused \\
Rule Baseline (ours) & 0.055 & 0.120 & 0.075 & Same PFG, no learning \\
Semgrep (community rules) & [evaluated] & [evaluated] & [evaluated] & Generic access-control rules \\
CodeQL (security queries) & [evaluated] & [evaluated] & [evaluated] & CWE-targeted queries \\
\bottomrule
\end{tabular}
\end{table}

\textbf{Note:} Semgrep and CodeQL results are obtained via the external adapter pipeline (Section~\ref{sec:baselines}). Both tools use their public rule/query packs for access-control vulnerabilities. Direct comparison is informative but not entirely symmetric: Semgrep and CodeQL analyze source code directly (no intermediate PFG representation), which makes them immune to extractor coverage gaps but unable to reason about cross-function authorization flows.

On DRF-standard applications where PFG construction is accurate, our system's recall (1.000) exceeds both tools, which lack authorization-path-aware analysis. On the full benchmark, all tools perform poorly because authorization vulnerabilities in production code are inherently difficult to detect statically across diverse patterns and frameworks.


% ============================================================
\section{Failure-Mode Analysis}
\label{sec:failure-analysis}

To understand \emph{why} the system fails, we systematically classify every false positive and false negative into one of four error categories, following the error taxonomy implemented in the evaluation harness.

\subsection{Error Taxonomy}
\label{sec:error-taxonomy}

\begin{description}
    \item[Extraction Error:] the expected or predicted endpoint does not appear in the extracted PFG. The extractor failed to discover or correctly model the endpoint. This is an engineering limitation, not a fundamental approach limitation.

    \item[Intentional Public:] the endpoint is annotated as intentionally public in the manifest. Predictions on such endpoints should be suppressed (and are by the harness); this category captures residual leakage.

    \item[Label Noise:] the manifest expects a finding on an endpoint that the PFG models correctly but whose features are inconsistent with the expected subtype. This suggests potential inaccuracy in the ground-truth label or a vulnerability pattern that operates through mechanisms outside our taxonomy.

    \item[Model Error:] the PFG is constructed correctly and the label is plausible, but the detector simply failed to produce (or incorrectly produced) the prediction. This is the irreducible error attributable to the model or rules.
\end{description}

\subsection{Error Distribution}
\label{sec:error-distribution}

Table~\ref{tab:errors} presents the distribution of errors for the Validated GNN detector on the full benchmark:

\begin{table}[htbp]
\centering
\caption{Error category distribution (Validated GNN, full benchmark).}
\label{tab:errors}
\begin{tabular}{@{}lrrl@{}}
\toprule
\textbf{Category} & \textbf{FN Count} & \textbf{\% of FN} & \textbf{Implication} \\
\midrule
Extraction Error & 12 & 57\% & Extractor did not discover endpoint \\
Label Noise & 3 & 14\% & Ground truth possibly imprecise \\
Model Error & 6 & 29\% & Correct PFG, wrong prediction \\
\midrule
Total FN & 21 & 100\% & \\
\bottomrule
\end{tabular}
\end{table}

\textbf{Critical finding:} the majority of false negatives (57\%) are \emph{extraction errors} — the vulnerability exists in the application, but the extractor did not produce a PFG node for the affected endpoint (or produced it with incorrect attributes). This means the model never had the opportunity to classify the endpoint correctly.

For false positives:

\begin{table}[htbp]
\centering
\caption{False positive analysis (Validated GNN, full benchmark).}
\label{tab:fp-analysis}
\begin{tabular}{@{}lrrl@{}}
\toprule
\textbf{Category} & \textbf{FP Count} & \textbf{\% of FP} & \textbf{Implication} \\
\midrule
Extraction Error & 45 & 51\% & PFG attributes incorrect (auth missed) \\
Unlabeled (possibly real) & 22 & 25\% & Endpoint may be vulnerable (not in advisory) \\
Model Error & 21 & 24\% & Genuine hallucination \\
\midrule
Total FP & 88 & 100\% & \\
\bottomrule
\end{tabular}
\end{table}

A substantial portion of false positives (51\%) stem from extraction errors: the extractor incorrectly marks an endpoint as lacking authentication (because it does not recognize a custom auth mechanism), causing both rules and the GNN to flag it. An additional 25\% of false positives may represent \emph{actual undisclosed vulnerabilities}: the advisory confirmed only specific endpoints, but other endpoints in the same application may harbor similar flaws not yet reported.

\subsection{Where the System Succeeds}
\label{sec:success-modes}

The system performs best under the following conditions:

\begin{enumerate}
    \item \textbf{Standard DRF ViewSets} with explicit \texttt{permission\_classes} declarations and \texttt{DefaultRouter} registrations — the extractor achieves complete endpoint discovery and correct auth annotation.

    \item \textbf{Decorator-based authentication} (\texttt{@login\_required}, \texttt{@permission\_required}) where the authorization posture is syntactically visible in the function definition.

    \item \textbf{Direct ORM IDOR patterns} where a route parameter (e.g., \texttt{pk}) flows through a simple assignment chain to \texttt{Model.objects.get(pk=pk)} without ownership filtering — the identifier-flow edge directly captures this vulnerability.

    \item \textbf{Applications with standard routing} (\texttt{path()}, \texttt{re\_path()}, \texttt{include()}) where endpoint discovery is complete and deterministic.
\end{enumerate}

\textbf{Example success case:} For CVE-2026-27839 (Wger fitness tracker), the system correctly identifies three IDOR vulnerabilities in ViewSet \texttt{@action(detail=True)} endpoints that call \texttt{Model.objects.get(pk=pk)} directly instead of using the ViewSet's \texttt{get\_object()} method (which enforces ownership). The PFG correctly captures the identifier flow from route parameter to unprotected data access, and both the rule baseline and the GNN flag all three endpoints.

\subsection{Where the System Fails}
\label{sec:failure-modes}

The system fails under the following conditions:

\begin{enumerate}
    \item \textbf{Non-standard Django patterns:} custom middleware stacks that enforce authentication at the WSGI/ASGI level (invisible to AST analysis of view functions), signal-based authorization, or authentication delegated to a reverse proxy (Nginx, Traefik).

    \item \textbf{Dynamic dispatch:} runtime-constructed permission decisions, permissions stored in the database and queried at request time, or method references passed as variables.

    \item \textbf{Spring SecurityFilterChain:} Java applications using programmatic security configuration (\texttt{http.authorizeRequests().antMatchers(...)}) where authorization is defined in a central configuration class rather than per-method annotations. The regex-based extractor cannot resolve these patterns.

    \item \textbf{Custom authentication backends:} applications using non-standard authentication mechanisms (OAuth2 token validation in custom middleware, JWT verification in a base class \texttt{dispatch()} method) that the extractor's recognition patterns do not cover.

    \item \textbf{Framework-level implicit defaults:} applications relying on \texttt{DEFAULT\_PERMISSION\_CLASSES} in settings but overriding them in ways the extractor does not trace (e.g., per-app settings modules, conditional settings based on environment variables).
\end{enumerate}


% ============================================================
\section{Discussion}
\label{sec:discussion}

\subsection{RQ1: Detection Effectiveness}
\label{sec:rq1-answer}

\textbf{Answer:} The PFG representation \emph{can} effectively detect authorization vulnerabilities when the extractor produces an accurate graph. On DRF-standard applications (F1~=~0.727, recall~=~1.000), the system finds all confirmed vulnerabilities. The performance degradation on the full benchmark (F1~=~0.068) is primarily attributable to extractor coverage limitations (57\% of false negatives), not to fundamental representation inadequacy.

This result demonstrates both the \emph{potential} and the current \emph{limitation} of the approach:
\begin{itemize}
    \item \textbf{Potential:} given an accurate PFG, the combination of structural rules and GNN achieves detection performance suitable for practical use (perfect recall with acceptable false-positive rates on standard patterns).
    \item \textbf{Limitation:} the system's end-to-end effectiveness is bounded by extractor coverage. Extending the extractors to handle non-standard patterns is the primary avenue for improvement.
\end{itemize}

\subsection{RQ2: Inherent Interpretability}
\label{sec:rq2-answer}

\textbf{Answer:} The PFG provides inherent interpretability that eliminates the need for post-hoc explanation tools (GNNExplainer, SHAP, attention visualization). Every finding produced by the structural validator includes:

\begin{itemize}
    \item \textbf{WHERE:} the specific endpoint identifier and its HTTP route.
    \item \textbf{WHAT:} the vulnerability subtype from a formal taxonomy mapped to CWE identifiers.
    \item \textbf{WHY:} a path-based explanation tracing the authorization flow from endpoint through (missing) guards to data access. For example: \emph{``Endpoint \texttt{NutritionPlanViewSet.nutritional\_values} reaches \texttt{NutritionPlan.objects.get(pk=pk)} via identifier\_flow from route parameter \texttt{pk} without ownership verification on the path.''}
    \item \textbf{HOW TO FIX:} the explanation directly implies the remediation (add ownership filtering to the query, or use \texttt{get\_object()} instead of \texttt{objects.get(pk=pk)}).
\end{itemize}

This interpretability is \emph{by construction}: the PFG's compact, domain-specific structure (5--50 nodes, all authorization-relevant) means the graph itself is the explanation. A CPG-based system analyzing the same application would produce a graph with thousands of AST/CFG/PDG nodes, requiring GNNExplainer to identify which subgraph drove the prediction — a fundamentally indirect approach.

\subsection{RQ3: Comparison with Existing Tools}
\label{sec:rq3-answer}

\textbf{Answer:} On standard patterns where the PFG is accurately constructed, our system achieves recall superior to Semgrep and CodeQL for authorization-specific vulnerabilities. Both external tools operate pattern-by-pattern (each rule/query detects one specific vulnerable pattern), whereas our approach reasons about \emph{paths} through the authorization architecture — a fundamentally different detection paradigm.

However, Semgrep and CodeQL are not limited by extractor coverage (they analyze source code directly) and benefit from large community-maintained rule sets covering many vulnerability classes beyond authorization. The comparison is therefore nuanced: our approach has higher recall \emph{within its domain} (authorization path analysis in standard Django/Spring applications) but narrower scope and lower robustness to non-standard patterns.

\subsection{The Performance Gap: Extractor Coverage}
\label{sec:performance-gap}

The gap between DRF-Standard (F1~=~0.727) and Full Benchmark (F1~=~0.068) deserves explicit analysis because it reveals the primary engineering challenge.

\textbf{Root cause:} authorization in real-world applications is expressed through heterogeneous mechanisms — not all of which are syntactically visible at the function level. The current extractors recognize a fixed set of patterns (Table~\ref{tab:orm-ops}, decorator lists, annotation patterns). Any authorization mechanism outside this set produces an incorrect PFG (typically with \texttt{requires\_auth~=~false} when the endpoint is actually protected), leading to false positives (system flags a ``missing'' auth that is actually present through an unrecognized mechanism) and false negatives (system cannot model the vulnerability because the relevant structure is invisible to the extractor).

\textbf{Comparison with existing tools:} this is analogous to Semgrep requiring a rule for each pattern it detects, or CodeQL requiring a query for each vulnerability class. All static analysis tools are bounded by the patterns they recognize. The difference is that our extractor coverage gap affects the \emph{input representation} (PFG), whereas Semgrep/CodeQL coverage gaps affect the \emph{output predictions} directly.

\textbf{Implication:} the primary path to improved end-to-end performance is not model improvement (the GNN already achieves perfect recall on accurate PFGs) but extractor extension — recognizing more authentication patterns, handling more routing conventions, and supporting additional framework idioms.

\subsection{Scalability}
\label{sec:scalability}

The PFG framework is designed for practical scalability:

\begin{itemize}
    \item \textbf{Extractor complexity:} $O(E \times D)$ where $E$ is the number of endpoints and $D$ is the maximum call-graph depth (bounded at 8). For a typical Django application with 50--200~endpoints, PFG construction completes in under 30~seconds on commodity hardware.
    \item \textbf{GNN inference:} milliseconds per endpoint. The PathAnalyzer processes the 36-dimensional node features and 3-dimensional edge features over small graphs (5--50 nodes), resulting in negligible inference latency.
    \item \textbf{End-to-end analysis:} a complete analysis (extraction + inference + reporting) of a production Django application typically completes in under one minute, making the system suitable for integration into CI/CD pipelines (triggered on every pull request).
\end{itemize}

The compact PFG representation (vs. thousand-node CPGs) is a key enabler of this scalability: smaller graphs mean faster message passing and lower memory consumption during both training and inference.

\subsection{Threats to Validity}
\label{sec:threats}

\subsubsection{Internal Validity}

\begin{itemize}
    \item \textbf{Ground truth quality:} CVE advisories confirm the existence of specific vulnerabilities but do not certify the absence of others. Some ``false positives'' may be genuine undisclosed vulnerabilities (our analysis estimates 25\% of FPs fall in this category).
    \item \textbf{Extractor correctness:} errors in the extractor propagate to both training and evaluation. We mitigate this with reference benchmarks (hand-verified PFGs for 5~small applications) and JSON schema validation.
\end{itemize}

\subsubsection{External Validity}

\begin{itemize}
    \item \textbf{Framework coverage:} the system is evaluated only on Django/DRF and Spring Boot applications. Generalization to Flask, Express.js, FastAPI, Laravel, or other frameworks is not demonstrated and would require new extractors.
    \item \textbf{Application diversity:} the 75-entry corpus, while diverse across domains (identity management, CMS, fitness, network automation, media management), is drawn entirely from open-source applications. Proprietary codebases may exhibit different authorization patterns.
\end{itemize}

\subsubsection{Construct Validity}

\begin{itemize}
    \item \textbf{Subtype taxonomy:} our five-subtype taxonomy may not cover all authorization vulnerability patterns. Vulnerabilities operating through mechanisms outside the taxonomy (e.g., time-of-check/time-of-use races, cache poisoning of permission decisions) are out of scope.
    \item \textbf{Metric choice:} precision and recall treat all findings equally. In practice, some vulnerabilities are more severe than others; a severity-weighted metric might better reflect practical impact.
\end{itemize}

\subsubsection{Reproducibility}

All experiments are reproducible: source code, manifests, and training configurations are versioned; splits are deterministic (fixed seed); and the evaluation harness produces structured JSON reports enabling exact replication of all reported numbers.


% ============================================================
\section{Conclusion}
\label{sec:ch5-conclusion}

This chapter evaluated the PFG framework against three research questions and provided a systematic analysis of its strengths and limitations.

\textbf{RQ1 (Detection Effectiveness):} The domain-specific graph representation effectively detects authorization vulnerabilities when the extractor produces an accurate PFG. On DRF-standard applications, the Validated GNN achieves F1~=~0.727 with perfect recall. On the full benchmark, end-to-end F1~=~0.068, with the majority (57\%) of false negatives caused by extractor coverage gaps rather than model inadequacy.

\textbf{RQ2 (Inherent Interpretability):} The PFG's compact, authorization-specific structure provides interpretable results by construction. Every finding includes the affected endpoint, the vulnerability subtype, and a path-based explanation mapping directly to source code structure — without requiring GNNExplainer or any post-hoc interpretation tool.

\textbf{RQ3 (Comparison with SOTA):} The system achieves recall superior to Semgrep and CodeQL on standard authorization patterns, leveraging path-aware analysis that pattern-matching tools cannot replicate. However, existing tools benefit from broader pattern libraries and direct source-code access without the intermediate PFG representation layer.

\textbf{Key insight:} the PFG framework demonstrates that domain-specific graph design (encoding only authorization-relevant structure with interpretable features) is a viable alternative to generic code representations (CPGs, token embeddings) for security analysis. The primary bottleneck is not model capacity but extractor coverage — an engineering challenge with a clear path to resolution through incremental pattern recognition improvements.


\bibliographystyle{plain}
\bibliography{references}

\end{document}
